\SourcePage{279}
\SGAchapter{XXVI}{约化群的抛物子群}{Sous-groupes paraboliques des groupes réductifs}
\begin{center}M. Demazure\end{center}
本讲研究\Term{reductive-group}{约化} $S$-群 $G$ 的\Term{parabolic-subgroup}{抛物子群}。主要结果是\Term{conjugacy}{共轭}定理（5.4）。主要工具是两个\Term{parabolic-subgroup}{抛物子群}的\Term{transversal-position}{横截位置}概念，第 4 节对此作系统研究。另一个重要事实是，\Term{parabolic-subgroup}{抛物子群} $P$ 的\Term{unipotent-radical}{幂幺根基} $\operatorname{rad}^{u}(P)$ \Term{decomposable}{可分解}为\Term{vector-group}{向量群}的逐次\Term{group-extension}{扩张}（2.1）。\footnote{编者注：由于 $P$ 的 \Term{levi-subgroup}{Levi 子群}构成一个以 $\operatorname{rad}^{u}(P)$ 为结构群的\Term{torsor}{挠子}（1.9），当 $S$ 半\Term{local-scheme}{局部}时，这推出 $P$ 有 \Term{levi-subgroup}{Levi 子群}，因而有\Term{maximal-torus}{极大环面}（2.4）。}
第 3 节研究与 $G$ 相关的若干\Term{scheme}{概形}；第 6 节研究 $G$ 的\Term{split-torus}{分裂子环面}\footnote{编者注：将“平凡\Term{torus}{环面}”一词改为“\Term{split-torus}{分裂环面}”。}及其与\Term{parabolic-subgroup}{抛物子群}的关系。
最后，第 7 节简要叙述如何在半\Term{local-scheme}{局部}\Term{base}{基底}上表述\Term{reductive-group}{约化群}的“相对理论”；域上的这种理论见 A. Borel 与 J. Tits 的文章 \textit{Groupes réductifs}, Publications Mathématiques de l'IHÉS, no.~27。下文以 [BT65] 引用该文；当\Term{base}{基底}是域时，读者还可在那里找到此处未曾涉及的其他结果。
\SGAsection{1}{回顾。Levi 子群}
\begin{SGAresult}{定义}{1.1}
设 $S$ 为\Term{scheme}{概形}，$G$ 为\Term{reductive-group}{约化} $S$-群，$P$ 为 $G$ 的子 $S$-群\Term{scheme}{概形}。若满足以下条件，则称 $P$ 为 $G$ 的\Term{parabolic-subgroup}{抛物子群}：
(i) $P$ 在 $S$ 上\Term{smooth}{光滑}；
(ii) 对每个 $s\in S$，商 $\bar s$-\Term{scheme}{概形} $G_{\bar s}/P_{\bar s}$ 是\Term{proper}{固有}的（即由 Bible，§6.4，定理 4（= [Ch05]，§6.5，定理 5），$P_{\bar s}$ 包含 $G_{\bar s}$ 的一个\Term{borel-subgroup}{Borel 子群}）。
\end{SGAresult}
\begin{SGAresult}{命题}{1.2}
（第 XXII 讲，5.8.5。）
设 $S$ 为\Term{scheme}{概形}，$G$ 为\Term{reductive-group}{约化} $S$-群，$P$ 为 $G$ 的\Term{parabolic-subgroup}{抛物子群}。则 $P$ 在 $G$ 中闭，\Term{fiber}{纤维}\Term{connected}{连通}，且
\[P=\operatorname{Norm}_G(P).\]
\SourcePage{280}
此外，\Term{quotient-sheaf}{商层} $G/P$ 可由一个在 $S$ 上\Term{smooth}{光滑}且\Term{projective}{射影}的 $S$-\Term{scheme}{概形}表示。
\end{SGAresult}
\begin{SGAresult}{命题}{1.3}
（第 XXII 讲，5.3.9 与 5.3.11。）
设 $S$ 为\Term{scheme}{概形}，$G$ 为\Term{reductive-group}{约化} $S$-群，$P,P'$ 为 $G$ 的两个\Term{parabolic-subgroup}{抛物子群}。下列条件\Term{equivalence}{等价}：
(i) $P,P'$ 在 $G$ 中关于\Term{etale}{平展}\Term{topology}{拓扑}（或 $(\mathrm{fpqc})$ \Term{topology}{拓扑}）\Term{local-scheme}{局部}\Term{conjugacy}{共轭}。
(ii) 对每个 $s\in S$，$P_s,P'_s$ 由 $G(\bar s)$ 的元素\Term{conjugacy}{共轭}。
(iii) $P$ 到 $P'$ 的严格\Term{transporter}{输运子} $\operatorname{Transp}^{\mathrm t}_G(P,P')$（对每个 $S'\to S$ 定义为
\[\operatorname{Transp}^{\mathrm t}_G(P,P')(S')=\{g\in G(S')\mid\operatorname{int}(g)P_{S'}=P'_{S'}\}\]）是 $G$ 的闭子\Term{scheme}{概形}，在 $S$ 上\Term{smooth}{光滑}且\Term{finite-presentation}{有限表现}，并且是右 $P$、左 $P'$ 的\Term{torsor}{主齐性丛}。
\end{SGAresult}
\begin{SGAresult}{命题}{1.4}
设 $S$ 为\Term{nonempty}{非空}\Term{scheme}{概形}，$(G,T,M,R)$ 为\Term{split-reductive-group}{分裂约化} $S$-群，$R'\subset R$。关于 $R'$ 的下列条件\Term{equivalence}{等价}：
(i) $\mathfrak g_{R'}=\mathfrak t\oplus\coprod_{\alpha\in R'}\mathfrak g^\alpha$ 是 $G$ 的一个包含 $T$ 的\Term{parabolic-subgroup}{抛物子群}的 \Term{lie-algebra}{Lie 代数}（该\Term{subgroup}{子群}必唯一，第 XXII 讲，5.3.5）。
(ii) $R'$ 为 $(R)$ 型（第 XXII 讲，5.4.2），且包含一个\Term{positive-root-system}{正根系}。
(iii) $R'$ 为 $R$ 的\Term{closed-subset}{闭子集}，并满足：若 $\alpha\in R-R'$，则 $-\alpha\in R'$（即 $R=R'\cup(-R')$）。
(iv) 存在一个\Term{simple-root-system}{单根系} $\Delta$ 及\Term{subset}{子集} $A\subset\Delta$，使 $R'$ 为所有\Term{positive-root}{正根}与所有能表示为 $A$ 中元素的线性组合的\Term{negative-root}{负根}之并。
(v) $R'$ 包含 $R$ 的一个\Term{simple-root-system}{单根系}；此外，若 $\Delta\subset R'$ 为 $R$ 的\Term{simple-root-system}{单根系}，并置 $A=(-R')\cap\Delta$，则 $R'$ 为所有\Term{positive-root}{正根}与所有能表示为 $A$ 中元素的线性组合的\Term{negative-root}{负根}之并。
\end{SGAresult}
\begin{proof}
由第 XXII 讲，5.4.5(ii)、5.5.1，有 (i)$\Leftrightarrow$(ii)；由第 XXI 讲，3.3.6 及第 XXII 讲，5.4.7，有 (iii)$\Rightarrow$(ii)。显然 (v)$\Rightarrow$(iv)$\Rightarrow$(iii)。由第 XXI 讲，3.3.6、3.3.10，有 (iii)$\Rightarrow$(v)。
故只需证明 (i) 推出 $R'$ 是 $R$ 的\Term{closed-subset}{闭子集}。该断言可在任意\Term{geometric-fibre}{几何纤维}上检验，故可假设 $S$ 为\Term{algebraically-closed-field}{代数闭域}的谱。
令 $P$ 为 $G$ 中包含 $T$、\Term{lie-algebra}{Lie 代数}为 $\mathfrak g_{R'}$ 的\Term{parabolic-subgroup}{抛物子群}。由于 $P$ 的\Term{borel-subgroup}{Borel 子群}就是包含在 $P$ 中的 $G$ 的\Term{borel-subgroup}{Borel 子群}，由 Bible，§12.3，定理 1 及第 XXII 讲，5.4.5(i)，若以 $U$ 表示 $P$ 的\Term{unipotent-radical}{幂幺根基}，则 $T\cdot U$ 是 $G$ 中包含 $T$、\Term{lie-algebra}{Lie 代数}为
\[\mathfrak g_{R''}=\mathfrak t\oplus\coprod_{\alpha\in R''}\mathfrak g^\alpha\]
的\Term{subgroup}{子群}。
\SourcePage{281}
这里 $R''$ 是所有包含于 $R'$ 的\Term{positive-root-system}{正根系}之交。特别地，$R''$ 闭且 $R''\cap(-R'')=\varnothing$。另一方面，$H=P/U$ 为\Term{reductive-group}{约化群}，$T$ 的\Term{canonical}{典范}像 $\overline T$ 是它的\Term{maximal-torus}{极大环面}（$T\to\overline T$ 是\Term{isomorphism}{同构}），并有 $T$-模\Term{isomorphism}{同构}，即 $M$ 型\Term{graded}{分次}\Term{vector-space}{向量空间}\Term{isomorphism}{同构}
\[\operatorname{Lie}(H)\simeq\mathfrak t\oplus\coprod_{\alpha\in R_s}\mathfrak g^\alpha,\]
其中 $R_s$ 是 $R''$ 在 $R'$ 中的补集。因此 $R_s$ 自然等同于 $H$ 关于 $\overline T$ 的\Term{root-set}{根集}，特别地 $R_s=-R_s$。这立即给出
\[R''=\{\alpha\in R'\mid-\alpha\notin R'\},\qquad R_s=\{\alpha\in R'\mid-\alpha\in R'\}.\]
现在证明 $R'$ 闭。设 $\alpha,\beta\in R'$ 且 $\alpha+\beta\in R$，须证 $\alpha+\beta\in R'$。若 $\alpha,\beta\in R''$，由 $R''$ 闭即得。若 $\alpha\in R_s,\beta\in R''$ 而 $\alpha+\beta\notin R'$，则 $\alpha+\beta\in-R''$，并由闭性得到 $-\alpha=-(\alpha+\beta)+\beta\in R''$；于是 $-\alpha\in R'$，因而 $-\alpha\in R_s$，与 $R_s\cap R''=\varnothing$ 矛盾。只剩 $\alpha,\beta\in R_s$ 的情形。若 $\alpha+\beta\notin R'$，则 $\alpha+\beta\in-R''$。但因 $\alpha+\beta\ne0$，\Term{root-system}{根系} $R_s$ 有一个包含 $\alpha,\beta$ 的\Term{positive-root-system}{正根系}，故 $H=P/U$ 有一个包含 $U_\alpha,U_\beta$ 的\Term{canonical}{典范}像的\Term{borel-subgroup}{Borel 子群}。它在 $P$ 中的\Term{inverse-image}{逆像}是一个包含 $U_\alpha,U_\beta,U$ 的\Term{borel-subgroup}{Borel 子群}，因而包含 $U_\alpha,U_\beta,U_{-(\alpha+\beta)}$，这是不可能的。
\end{proof}
\begin{SGAresult}{推论}{1.5}
\Term{reductive-group}{约化群}的\Term{parabolic-subgroup}{抛物子群}为 $(\mathrm{RC})$ 型（第 XXII 讲，5.11.1）。
\end{SGAresult}
\begin{SGAresult}{命题}{1.6}
（第 XXII 讲，5.11.4。）
设 $S$ 为\Term{scheme}{概形}，$G$ 为\Term{reductive-group}{约化} $S$-群，$P$ 为 $G$ 的一个 $(\mathrm{RC})$ 型\Term{subgroup}{子群}。\footnote{编者注：依照第 XXII 讲，5.11.4，将“\Term{parabolic-subgroup}{抛物子群}”改成“$(\mathrm{RC})$ 型\Term{subgroup}{子群}”，因为下文（4.5.1、6.17）需要在两个\Term{parabolic-subgroup}{抛物子群} $P,P'$ 的交 $P\cap P'$ 为 $(\mathrm{RC})$ 型时对其应用此命题。}
(i) $P$ 有一个最大的\Term{normal-subgroup}{正规子群}\Term{scheme}{概形}，在 $S$ 上\Term{smooth}{光滑}且\Term{finite-presentation}{有限表现}，\Term{geometric-fibre}{几何纤维}\Term{connected}{连通}且幂幺。它是 $P$ 的\Term{characteristic-subgroup}{特征子群}，称为 $P$ 的\Term{unipotent-radical}{幂幺根基}，记作 $\operatorname{rad}^{u}(P)$。\Term{quotient-sheaf}{商层} $P/\operatorname{rad}^{u}(P)$ 可由\Term{reductive-group}{约化} $S$-群表示。
(ii) 若 $T$ 为 $P$ 的\Term{maximal-torus}{极大环面}，则 $P$ 有一个包含 $T$ 的\Term{reductive-group}{约化子群} $L$，满足：
(a) $P$ 中包含 $T$ 的每个\Term{reductive-group}{约化子群}均包含在 $L$ 中；
(b) $P$ 为\Term{semidirect-product}{半直积} $L\cdot\operatorname{rad}^{u}(P)$，即\Term{canonical}{典范}\Term{morphism}{态射} $L\to P/\operatorname{rad}^{u}(P)$ 为\Term{isomorphism}{同构}。
此外，$L$ 是 $P$ 中包含 $T$ 且满足 (b) 的唯一\Term{subgroup}{子群}（或满足 (a) 的唯一\Term{reductive-group}{约化子群}）。最后，
\[\operatorname{Norm}_P(L)=L,\qquad\operatorname{Norm}_P(T)=\operatorname{Norm}_L(T).\]
\end{SGAresult}
\SourcePage{282}
\textbf{1.7.} $P$ 中满足上述条件 (b) 的\Term{subgroup}{子群} $L$ 称为 $P$ 的\Term{levi-subgroup}{Levi 子群}。它是 $P$ 的极大\Term{reductive-group}{约化子群}：因它与 $P/\operatorname{rad}^{u}(P)$ \Term{isomorphism}{同构}，所以\Term{reduced}{约化}；为证极大性，设 $L'$ 为 $P$ 中包含 $L$ 的\Term{reductive-group}{约化子群}。证明 $L'=L$ 可关于 $(\mathrm{fpqc})$ \Term{topology}{拓扑}\Term{local-scheme}{局部}进行，故可假设 $L$ 有\Term{maximal-torus}{极大环面} $T$，于是归结于 1.6(ii)。
若 $L,L'$ 是 $P$ 的两个\Term{levi-subgroup}{Levi 子群}，则它们在 $P$ 中关于 $(\mathrm{fpqc})$ \Term{topology}{拓扑}\Term{local-scheme}{局部}\Term{conjugacy}{共轭}。事实上，可关于此\Term{topology}{拓扑}\Term{local-scheme}{局部}假设 $L$（或 $L'$）有\Term{maximal-torus}{极大环面} $T$（或 $T'$）；因 $T,T'$ 在 $P$ 中关于此\Term{topology}{拓扑}\Term{local-scheme}{局部}\Term{conjugacy}{共轭}，可假设 $T=T'$，再由 1.6(ii) 得 $L=L'$。由 $P=L\cdot\operatorname{rad}^{u}(P)$ 及 $\operatorname{Norm}_P(L)=L$ 立即得到：
\begin{SGAresult}{推论}{1.8}
设 $P$ 是\Term{reductive-group}{约化} $S$-群 $G$ 的 $(\mathrm{RC})$ 型\Term{subgroup}{子群}。\footnote{编者注：参见脚注 (3)。}若 $L,L'$ 为 $P$ 的两个\Term{levi-subgroup}{Levi 子群}，则存在唯一 $u\in\operatorname{rad}^{u}(P)(S)$ 使 $\operatorname{int}(u)L=L'$。
\end{SGAresult}
以 $\operatorname{Lev}(P)$ 表示 $P$ 的\Term{levi-subgroup}{Levi 子群}\Term{functor}{函子}：对 $S'\to S$，$\operatorname{Lev}(P)(S')$ 是 $P_{S'}$ 的\Term{levi-subgroup}{Levi 子群}集。由 1.8 得：
\begin{SGAresult}{推论}{1.9}
设 $P$ 是\Term{reductive-group}{约化} $S$-群 $G$ 的 $(\mathrm{RC})$ 型\Term{subgroup}{子群}（参见脚注 (3)）。则 $\operatorname{Lev}(P)$ 是以 $S$-群 $\operatorname{rad}^{u}(P)$ 为结构群的\Term{torsor}{主齐性丛}；特别地，可由在 $S$ 上\Term{smooth}{光滑}且\Term{affine}{仿射}、\Term{geometric-fibre}{几何纤维}整的 $S$-\Term{scheme}{概形}表示。
\end{SGAresult}
由 1.6 立即得到：
\begin{SGAresult}{推论}{1.10}
设 $P$ 是\Term{reductive-group}{约化} $S$-群 $G$ 的 $(\mathrm{RC})$ 型\Term{subgroup}{子群}（参见脚注 (3)）。$P$ 的\Term{maximal-torus}{极大环面}\Term{functor}{函子} $\operatorname{Tor}(P)$ 可由\Term{smooth}{光滑}\Term{affine}{仿射} $S$-\Term{scheme}{概形}表示。“$L\supset T$”这一关系定义\Term{morphism}{态射}
\[\operatorname{Tor}(P)\longrightarrow\operatorname{Lev}(P),\]
它在 $L\in\operatorname{Lev}(P)(S)$ 上方的\Term{fiber}{纤维}等同于 $\operatorname{Tor}(L)$（第 XXII 讲，5.8.3）。
\end{SGAresult}
1.10 的第一个断言由另两个断言及第 XXII 讲，5.8.3 推出。
\begin{SGAresult}{定义}{1.11}
设 $S$ 为\Term{nonempty}{非空}\Term{scheme}{概形}，$G$ 为\Term{reductive-group}{约化} $S$-群，$P$ 为 $G$ 的\Term{parabolic-subgroup}{抛物子群}，$E=(T,M,R,\Delta,(X_\alpha)_{\alpha\in\Delta})$ 为 $G$ 的\Term{pinning}{钉扎}。若 $P\supset T$ 且 $P$ 的 \Term{lie-algebra}{Lie 代数}形如 $\mathfrak t\oplus\coprod_{\alpha\in R'}\mathfrak g^\alpha$，其中 $R'\subset R$ 包含 $R^+$，则称 $E$ 适合于 $P$，或称 $E$ 为对 $(G,P)$ 的\Term{pinning}{钉扎}。
\end{SGAresult}
特别地，若 $T\subset B$ 是此\Term{pinning}{钉扎}定义的 Killing 对，则 $T\subset B\subset P$。在以上条件下，记 $\Delta(P)=\Delta\cap(-R')$；由第 XXII 讲，5.4.3，
\[\alpha\in\Delta(P)\ \Longleftrightarrow\ \alpha\in\Delta\text{ 且 }U_{-\alpha}\subset P\ \Longleftrightarrow\ \alpha\in\Delta\text{ 且 }U_{-\alpha}\cap P\ne e.\]
\SourcePage{283}
由 1.4(v) 及第 XXII 讲，5.11.3、5.10.6 立即得到：
\begin{SGAresult}{命题}{1.12}
设 $S$ 为\Term{scheme}{概形}，$G$ 为\Term{reductive-group}{约化} $S$-群，$P$ 为 $G$ 的\Term{parabolic-subgroup}{抛物子群}，$(T,M,R,\Delta,(X_\alpha)_{\alpha\in\Delta})$ 为适合于 $P$ 的 $G$ 的\Term{pinning}{钉扎}，$\Delta(P)$ 为以上定义的 $\Delta$ 的\Term{subset}{子集}。
(i) $P$ 的\Term{unipotent-radical}{幂幺根基} $\operatorname{rad}^{u}(P)$ 正是
\[U_{R''}=\prod_{\alpha\in R''}U_\alpha,\]
其中 $R''$ 为那些在关于 $\Delta$ 的\Term{resolution}{分解}中，至少有一个 $\Delta-\Delta(P)$ 的元素系数非零的\Term{positive-root}{正根}所成的\Term{set}{集合}。\footnote{编者注：将 $\Delta(P)$ 纠正为 $\Delta-\Delta(P)$。}
(ii) $P$ 中包含 $T$ 的唯一\Term{levi-subgroup}{Levi 子群} $L$ 正是
\[Z_{\Delta(P)}=\operatorname{Centr}_G(T_{\Delta(P)}),\]
其中 $T_{\Delta(P)}$ 是 $\bigcap_{\alpha\in\Delta(P)}\operatorname{Ker}(\alpha)$ 的\Term{maximal-torus}{极大环面}；此外 $T_{\Delta(P)}=\operatorname{rad}(L)$。
\end{SGAresult}
\begin{SGAresult}{推论}{1.13}
\Term{reductive-group}{约化群} $G$ 的\Term{parabolic-subgroup}{抛物子群} $P$ 的每个\Term{levi-subgroup}{Levi 子群} $L$ 都是 $G$ 的\Term{critical-subgroup}{临界子群}，即满足（第 XXII 讲，5.10.4）
\[L=\operatorname{Centr}_G(\operatorname{rad}(L)).\]
\end{SGAresult}
这由 1.12 及包含在 1.4 和第 XXII 讲，5.4.1 中的下述引理立即推出：
\begin{SGAresult}{引理}{1.14}
关于\Term{etale}{平展}\Term{topology}{拓扑}\Term{local-scheme}{局部}地，每一对 $(G,P)$，其中 $P$ 为\Term{reductive-group}{约化群} $G$ 的\Term{parabolic-subgroup}{抛物子群}，都可赋予\Term{pinning}{钉扎}（1.11）。
\end{SGAresult}
还须指出：
\begin{SGAresult}{命题}{1.15}
设 $S$ 为\Term{scheme}{概形}，$G$ 为\Term{reductive-group}{约化} $S$-群，$P$ 为 $G$ 的\Term{parabolic-subgroup}{抛物子群}，$E,E'$ 为适合于 $P$ 的两个 $G$ 的\Term{pinning}{钉扎}。将 $E$ 变为 $E'$ 的 $G$ 在 $S$ 上的唯一\Term{inner-automorphism}{内自同构}（第 XXIV 讲，1.5）由下列\Term{morphism}{态射}来自 $P$：
\[P\longrightarrow P/\operatorname{Centr}(P)=P/\operatorname{Centr}(G)\longrightarrow G/\operatorname{Centr}(G).\]
\end{SGAresult}
事实上，只需仿照第 XXIV 讲，1.5 的论证，并使用：
\begin{SGAresult}{引理}{1.16}
（第 XXII 讲，5.3.14 与 5.2.6。）
\Term{reductive-group}{约化} $S$-群 $G$ 的\Term{parabolic-subgroup}{抛物子群} $P$ 的\Term{maximal-torus}{极大环面}（或\Term{borel-subgroup}{Borel 子群}，或 Killing 对）在 $P$ 中关于\Term{etale}{平展}\Term{topology}{拓扑}\Term{local-scheme}{局部}\Term{conjugacy}{共轭}。
\end{SGAresult}
\begin{SGAresult}{命题}{1.17}
设 $S$ 为\Term{scheme}{概形}，$G$ 为\Term{reductive-group}{约化} $S$-群，$P,P'$ 为 $G$ 的两个\Term{parabolic-subgroup}{抛物子群}，$B$ 为包含于 $P,P'$ 的\Term{borel-subgroup}{Borel 子群}。若 $P,P'$ 在 $G$ 中关于\Term{etale}{平展}\Term{topology}{拓扑}\Term{local-scheme}{局部}\Term{conjugacy}{共轭}，则 $P=P'$。
\end{SGAresult}
\begin{proof}
可假设存在 $g\in G(S)$ 使 $\operatorname{int}(g)P=P'$。于是 $B$ 与 $\operatorname{int}(g)^{-1}B$ 是 $P$ 的两个\Term{borel-subgroup}{Borel 子群}。
\SourcePage{284}
\Term{group-extension}{扩张} $S$ 后，由 1.16 可假设有 $p\in P(S)$ 使 $\operatorname{int}(p)\operatorname{int}(g^{-1})B=B$。于是
\[pg^{-1}\in\operatorname{Norm}_G(B)(S)=B(S),\]
故 $g\in B(S)\cdot p\subset P(S)$，从而 $P'=\operatorname{int}(g)P=P$。
\end{proof}
\textbf{注记 1.18。}若 $G$ 的两个\Term{parabolic-subgroup}{抛物子群} $P,P'$ 包含同一个\Term{borel-subgroup}{Borel 子群}，则 $P\cap P'$ 仍是 $G$ 的\Term{parabolic-subgroup}{抛物子群}。事实上，它沿\Term{unit-section}{单位截面}\Term{smooth}{光滑}（第 XXII 讲，5.4.5），且包含一个\Term{borel-subgroup}{Borel 子群}。
\begin{SGAresult}{命题}{1.19}
设 $S$ 为\Term{scheme}{概形}，$G$ 为\Term{reductive-group}{约化} $S$-群，$G'$ 为其\Term{derived-group}{导群}（第 XXII 讲，6.2.1）。
(i) 映射
\[P\longmapsto P'=P\cap G',\qquad P'\longmapsto P'\cdot\operatorname{rad}(G)=\operatorname{Norm}_G(P')\]
给出 $G$ 的\Term{parabolic-subgroup}{抛物子群}集与 $G'$ 的\Term{parabolic-subgroup}{抛物子群}集之间互逆的\Term{bijection}{双射}。且 $\operatorname{rad}^{u}(P)=\operatorname{rad}^{u}(P')$。
(ii) 设 $P$ 为 $G$ 的\Term{parabolic-subgroup}{抛物子群}，$P'=P\cap G'$。映射
\[\begin{aligned}L&\longmapsto L'=L\cap G'=L\cap P',\\L'&\longmapsto L'\cdot\operatorname{rad}(G)=\operatorname{Centr}_G(\operatorname{rad}(L'))\end{aligned}\]
给出 $P$ 的\Term{levi-subgroup}{Levi 子群}集与 $P'$ 的\Term{levi-subgroup}{Levi 子群}集之间互逆的\Term{bijection}{双射}。此外 $\operatorname{rad}(L')=(\operatorname{rad}(L)\cap G')^0$。
\end{SGAresult}
证明没有困难（例如可归约到\Term{splitting}{分裂}情形），留给读者；下述直接的命题也一样。
\begin{SGAresult}{命题}{1.20}
设 $S$ 为\Term{scheme}{概形}，$G$ 为\Term{reductive-group}{约化} $S$-群，$P$ 为 $G$ 的\Term{parabolic-subgroup}{抛物子群}，$L$ 为 $P$ 的\Term{levi-subgroup}{Levi 子群}。映射
\[Q\longmapsto Q\cap L=Q',\qquad Q'\longmapsto Q'\cdot\operatorname{rad}^{u}(P)\]
给出 $G$ 中包含于 $P$ 的\Term{parabolic-subgroup}{抛物子群}集与 $L$ 的\Term{parabolic-subgroup}{抛物子群}集之间互逆的\Term{bijection}{双射}。此外，$Q'$ 的\Term{levi-subgroup}{Levi 子群}正是 $Q$ 中包含于 $L$ 的\Term{levi-subgroup}{Levi 子群}。
\end{SGAresult}
可如下补充 1.6：
\begin{SGAresult}{命题}{1.21}
设 $S$ 为\Term{scheme}{概形}，$G$ 为\Term{reductive-group}{约化} $S$-群，$P$ 为 $G$ 的\Term{parabolic-subgroup}{抛物子群}。
(i) $P$ 有一个最大的\Term{normal-subgroup}{正规子群}，在 $S$ 上\Term{smooth}{光滑}且\Term{finite-presentation}{有限表现}，\Term{geometric-fibre}{几何纤维}\Term{connected}{连通}且\Term{solvable}{可解}。它是 $P$ 的\Term{characteristic-subgroup}{特征子群}，称为 $P$ 的\Term{radical}{根基}，记作 $\operatorname{rad}(P)$。\Term{quotient-sheaf}{商层} $P/\operatorname{rad}(P)$ 可由\Term{semisimple}{半单} $S$-群表示。
\SourcePage{285}
(ii) 若 $L$ 为 $P$ 的\Term{levi-subgroup}{Levi 子群}，则 $\operatorname{rad}(P)$ 为 $\operatorname{rad}^{u}(P)$ 与 $\operatorname{rad}(L)$ 的\Term{semidirect-product}{半直积}；有 $\operatorname{rad}(L)=L\cap\operatorname{rad}(P)$，从而 $L=\operatorname{Centr}_G(L\cap\operatorname{rad}(P))$，且 $P/\operatorname{rad}(P)\simeq L/\operatorname{rad}(L)$。
\end{SGAresult}
\begin{proof}
断言 (i) 是\Term{local-scheme}{局部}的，故可假设 $G$ 有\Term{levi-subgroup}{Levi 子群} $L$，于是只须证明 $R=\operatorname{rad}^{u}(P)\cdot\operatorname{rad}(L)$ 有 (i) 所述性质，这是直接的。对于 (ii)，只须再证 $\operatorname{rad}(L)=L\cap\operatorname{rad}(P)$；由于 $L$ 为一个\Term{torus}{环面}的\Term{centralizer}{中心化子}，$L\cap\operatorname{rad}(P)$ \Term{smooth}{光滑}且\Term{fiber}{纤维}\Term{connected}{连通}，所求立即成立。
\end{proof}
\SGAsection{2}{抛物子群的幂幺根基的结构}
\begin{SGAresult}{命题}{2.1}
设 $S$ 为\Term{scheme}{概形}，$G$ 为\Term{reductive-group}{约化} $S$-群，$P$ 为 $G$ 的\Term{parabolic-subgroup}{抛物子群}，$\operatorname{rad}^{u}(P)$ 为其\Term{unipotent-radical}{幂幺根基}。存在 $\operatorname{rad}^{u}(P)$ 的\Term{subgroup}{子群}\Term{scheme}{概形}列
\[\operatorname{rad}^{u}(P)=U_0\supset U_1\supset U_2\supset\cdots\supset U_n\supset\cdots\]
具有下列性质：
(i) 每个 $U_i$ \Term{smooth}{光滑}、\Term{fiber}{纤维}\Term{connected}{连通}，并是 $P$ 的闭\Term{characteristic-subgroup}{特征子群}。$U_i$ 的\Term{section}{截面}与 $U_j$ 的\Term{section}{截面}的交换子是 $U_{i+j+1}$ 的\Term{section}{截面}（在变动的 $S'\to S$ 上）。
(ii) 对每个 $i\ge0$，存在\Term{locally-free}{局部自由} $\mathcal O_S$-模 $\mathcal E_i$ 及 $S$-群\Term{sheaf}{层}的\Term{isomorphism}{同构}
\[U_i/U_{i+1}\xrightarrow{\sim}\mathbf W(\mathcal E_i).\]
此外，$P$ 的\Term{automorphism}{自同构}（在变动的 $S'\to S$ 上）线性地\Term{action}{作用}于 $\mathbf W(\mathcal E_i)$。
(iii) 对每个 $s\in S$，当 $n>\dim(\operatorname{rad}^{u}(P)_s)$ 时有 $U_{n,s}=e$。
\end{SGAresult}
\textbf{2.1.1.} 先假设对 $(G,P)$ 可\Term{pinning}{钉扎}。取适合于 $P$ 的 $G$ 的\Term{pinning}{钉扎} $(T,M,R,\Delta,\ldots)$，令 $\Delta(P)$ 为 $P$ 定义的 $\Delta$ 的\Term{subset}{子集}。设 $\alpha_1,\ldots,\alpha_p$ 为 $\Delta(P)$ 的元素，$\beta_1,\ldots,\beta_q$ 为 $\Delta-\Delta(P)$ 的元素。每个根 $\gamma\in R$ 唯一写作
\[\gamma=a_1\alpha_1+\cdots+a_p\alpha_p+b_1\beta_1+\cdots+b_q\beta_q.\]
置
\[a(\gamma)=b_1+\cdots+b_q.\]
\footnote{编者注：将 $a_1+\cdots+a_p$ 纠正为 $b_1+\cdots+b_q$。}
从定义立即得到以下性质（参见 1.12）：
\[\begin{aligned}
{}&\text{(i)}&&U_\gamma\subset P\ \Longleftrightarrow\ a(\gamma)\ge0;\\
&\text{(ii)}&&U_\gamma\subset\operatorname{rad}^{u}(P)\ \Longleftrightarrow\ a(\gamma)>0;\\
&\text{(iii)}&&a(n\gamma+m\gamma')=na(\gamma)+ma(\gamma')\quad(n,m\in\mathbf Z).
\end{aligned}\]
\SourcePage{286}
对 $i>0$，令 $R_i$ 为满足 $a(\gamma)>i$ 的根 $\gamma\in R$ 所成的\Term{set}{集合}。每个 $R_i$ 是\Term{closed-root-set}{闭根集}，且 $R_i\cap(-R_i)=\varnothing$。考虑 $S$-群（第 XXII 讲，5.6.5）
\[U_i=U_{R_i}=\prod_{\gamma\in R_i}U_\gamma.\]
它是 $G$ 的\Term{closed-subgroup}{闭子群}\Term{scheme}{概形}，在 $S$ 上\Term{smooth}{光滑}且\Term{fiber}{纤维}\Term{connected}{连通}。
设 $\alpha,\beta\in R$，考虑第 XXII 讲，5.5.2 的交换关系
\[p_\alpha(x)p_\beta(y)p_\alpha(-x)=p_\beta(y)\prod_{n,m\in\mathbf N^*}p_{n\alpha+m\beta}(C_{n,m,\alpha,\beta}x^ny^m),\]
其中每个 $p_\gamma$ 是\Term{vector-group}{向量群}\Term{isomorphism}{同构} $\mathbf G_{a,S}\xrightarrow{\sim}U_\gamma$。先注意，若 $a(\alpha)>i$、$a(\beta)>j$，则当 $n,m>0$ 时
\[a(n\alpha+m\beta)=na(\alpha)+ma(\beta)\ge n(i+1)+m(j+1)>i+j+1.\]
故 $U_\alpha$ 的\Term{section}{截面}与 $U_\beta$ 的\Term{section}{截面}的交换子是 $U_{i+j+1}$ 的\Term{section}{截面}（在变动的 $S'\to S$ 上），这确实推出 $(U_i,U_j)\subset U_{i+j+1}$。故对每个 $i\ge0$，商 $U_i/U_{i+1}$ 交换，且自然等同于
\[U_i/U_{i+1}\simeq\prod_{a(\gamma)=i+1}U_\gamma\simeq\mathbf W(\mathcal E_i),\]
其中 $\mathcal E_i$ 是所有满足 $a(\gamma)=i+1$ 的 $\mathfrak g^\gamma$ 的\Term{direct-sum}{直和}。
回到上述交换公式，假设 $a(\alpha)\ge0$、$a(\beta)>i$。若 $n,m\in\mathbf N^*$，则或有 $a(n\alpha+m\beta)>i+1$，或有 $a(n\alpha+m\beta)=i+1$，后一情形必有 $m=1$。这首先说明 $\operatorname{int}(p_\alpha(x))$ 保持 $U_i$（从而也保持 $U_{i+1}$）；其次说明，在 $\operatorname{int}(p_\alpha(x))p_\beta(y)$ 的表达式中，模 $U_{i+1}$ 只出现形如 $p_{n\alpha+\beta}(C_{n,1,\alpha,\beta}x^ny)$ 的项，故关于 $y$ 线性。因此，由 $U_\alpha$ 的\Term{section}{截面}定义的\Term{inner-automorphism}{内自同构}在等同于 $\mathbf W(\mathcal E_i)$ 的商 $U_i/U_{i+1}$ 上线性\Term{action}{作用}。对 $T$ 的\Term{section}{截面}定义的\Term{inner-automorphism}{内自同构}，这显然也成立；又 $P$ 由 $T$ 和满足 $a(\alpha)\ge0$ 的 $U_\alpha$ 生成，故得到：
(i) 每个 $U_i$ 在 $P$ 中\Term{normal}{正规}；
(ii) $P$ 的\Term{section}{截面}定义的\Term{inner-automorphism}{内自同构}线性地\Term{action}{作用}于 $U_i/U_{i+1}\simeq\mathbf W(\mathcal E_i)$。
\textbf{2.1.2.} 再取适合于 $P$ 的一个新\Term{pinning}{钉扎} $(T',M',R',\Delta',\ldots)$。由 1.15，有一个来自 $P$ 的 $G$ 的\Term{inner-automorphism}{内自同构}将旧\Term{pinning}{钉扎}变成新\Term{pinning}{钉扎}。\Term{group-extension}{扩张} $S$ 后，可假设它为 $\operatorname{int}(p)$，$p\in P(S)$。用新\Term{pinning}{钉扎}重复以上构造，显然所得群 $U'_i$ 及\Term{isomorphism}{同构} $U'_i/U'_{i+1}\simeq\mathbf W(\mathcal E'_i)$，由 $U_i$ 及 $U_i/U_{i+1}\simeq\mathbf W(\mathcal E_i)$ 经 $\operatorname{int}(p)$ \Term{transport-of-structure}{结构搬运}而得。由上述 (i)、(ii)，
\SourcePage{287}
$U'_i=U_i$，且在 $U_i/U_{i+1}=U'_i/U'_{i+1}$ 上构造的两个向量结构相同。
这说明群 $U_i$ 及商 $U_i/U_{i+1}$ 上的向量结构不依赖于所取\Term{pinning}{钉扎}（特别地，易见它们在 $P$ 的一切\Term{automorphism}{自同构}下不变）。
因此已在对 $(G,P)$ 可\Term{pinning}{钉扎}时证明了命题（(iii) 显然，因为由第 XXI 讲，3.1.2，\Term{set}{集合} $\{a(\gamma)\mid\gamma\in R\}$ 为 $\mathbf Z$ 的区间，故仅当 $U_{i,s}=e$ 时才可能有 $\dim U_{i,s}=\dim U_{i+1,s}$）。
\textbf{2.1.3.} 一般情形下，有 $(\mathrm{fpqc})$ \Term{topology}{拓扑}的\Term{covering-family}{覆盖族} $\{S_j\to S\}$，使每对 $(G_{S_j},P_{S_j})$ 可\Term{pinning}{钉扎}（1.14）。由以上讨论，$\operatorname{rad}^{u}(P_{S_j})_i$ 上有与商的向量结构相容的\Term{descent}{下降}数据，由闭子\Term{scheme}{概形}（或\Term{locally-free}{局部自由}模）的\Term{descent}{下降}便得结论。\footnote{编者注：参见 SGA 1，第 VIII 讲，1.3、1.9、1.10。}
\begin{SGAresult}{推论}{2.2}
设 $S$ 为\Term{affine}{仿射}\Term{scheme}{概形}，$G$ 为\Term{reductive-group}{约化} $S$-群，$P$ 为 $G$ 的\Term{parabolic-subgroup}{抛物子群}。有
\[\mathrm H^1(S,\operatorname{rad}^{u}(P))=0,\]
即以 $\operatorname{rad}^{u}(P)$ 为结构群的每个\Term{torsor}{主齐性丛}均平凡。
\end{SGAresult}
\begin{proof}
$S$ \Term{decomposable}{可分解}为子\Term{scheme}{概形}之和，使 $\operatorname{rad}^{u}(P)$ 在每个子\Term{scheme}{概形}上\Term{relative-dimension}{相对维数}恒定。由 (iii)，可假设存在 $n$ 使 $U_n=e$。由 TDTE I，B，1.1（或 SGA 1，第 XI 讲，5.1），有 $\mathrm H^1(S,U_i/U_{i+1})=\mathrm H^1(S,\mathbf W(\mathcal E_i))=0$，结论立即成立。
\end{proof}
\begin{SGAresult}{推论}{2.3}
在以上条件下，$P$ 有\Term{levi-subgroup}{Levi 子群} $L$。若 $L$ 是 $P$ 的\Term{levi-subgroup}{Levi 子群}，则\Term{canonical}{典范}映射
\[\mathrm H^1(S,L)\longrightarrow\mathrm H^1(S,P)\]
为\Term{bijection}{双射}（$\mathrm H^1(S,\ )$ 的定义参见第 XXIV 讲引言）。
\end{SGAresult}
\begin{proof}
第一个断言由 2.2 及 1.9 推出。因 $P$ 为\Term{semidirect-product}{半直积} $L\cdot\operatorname{rad}^{u}(P)$，\Term{canonical}{典范}映射 $\mathrm H^1(S,L)\to\mathrm H^1(S,P)$ \Term{surjectivity}{满射}。为证\Term{injectivity}{单射}，只需说明对每个以 $L$ 为结构群的\Term{torsor}{主齐性丛} $Q$，有 $\mathrm H^1(S,\operatorname{rad}^{u}(P)_Q)=0$，其中下标 $Q$ 表示由 $L$-丛 $Q$ 作\Term{twisting}{扭曲}。可用两种方式证明：利用 $U_i/U_{i+1}$ 上的向量结构在 $L$ 下不变，重复 2.2 的证明；或者注意 $\operatorname{rad}^{u}(P)_Q$ 等同于 $G_Q$ 的\Term{parabolic-subgroup}{抛物子群} $P_Q$ 的\Term{unipotent-radical}{幂幺根基}，然后对 $P_Q$ 应用 2.2。
\end{proof}
\begin{SGAresult}{推论}{2.4}
设 $S$ 为半\Term{local-scheme}{局部}\Term{scheme}{概形}，$G$ 为\Term{reductive-group}{约化} $S$-群，$P$ 为 $G$ 的\Term{parabolic-subgroup}{抛物子群}。存在 $G$ 的\Term{maximal-torus}{极大环面} $T$ 包含在 $P$ 中。
\end{SGAresult}
事实上，由 2.3，$P$ 有\Term{levi-subgroup}{Levi 子群} $L$，只须证明 $L$ 有\Term{maximal-torus}{极大环面}，而这由第 XIV 讲，3.20 推出。
\SourcePage{288}
\begin{SGAresult}{推论}{2.5}
设 $S$ 为\Term{affine}{仿射}\Term{scheme}{概形}，$G$ 为\Term{reductive-group}{约化} $S$-群，$P$ 为 $G$ 的\Term{parabolic-subgroup}{抛物子群}。存在\Term{locally-free}{局部自由} $\mathcal O_S$-模 $\mathcal E$，使 $\operatorname{rad}^{u}(P)$ 作为 $S$-\Term{scheme}{概形}\Term{isomorphism}{同构}于 $\mathbf W(\mathcal E)$。
\end{SGAresult}
\begin{proof}
对 $i$ 归纳证明有 $S$-\Term{scheme}{概形}\Term{isomorphism}{同构}
\[\operatorname{rad}^{u}(P)/U_i\simeq\mathbf W(\mathcal E_0\oplus\cdots\oplus\mathcal E_{i-1}).\]
$i=0$ 时显然。设 $i>0$，则 $\operatorname{rad}^{u}(P)/U_i$ 为以 $\operatorname{rad}^{u}(P)/U_{i-1}$ 为\Term{base}{基底}、以群
\[(U_{i-1}/U_i)_{\operatorname{rad}^{u}(P)/U_{i-1}}\simeq\mathbf W(\mathcal E_{i-1}\otimes\mathcal O_{\operatorname{rad}^{u}(P)/U_{i-1}})\]
为结构群的\Term{torsor}{主齐性丛}。\Term{base}{基底}\Term{affine}{仿射}（例如由归纳假设），故此丛平凡（TDTE I 或 SGA 1，第 XI 讲，上引处），因而有 $S$-\Term{scheme}{概形}\Term{isomorphism}{同构} $\operatorname{rad}^{u}(P)/U_i\simeq(\operatorname{rad}^{u}(P)/U_{i-1})\times_S(U_{i-1}/U_i)$，证明完成。
\end{proof}
\begin{SGAresult}{推论}{2.6}
设 $S$ 为半\Term{local-scheme}{局部}\Term{scheme}{概形}，$\{s_i\}$ 为其\Term{closed-point}{闭点}集，$G$ 为\Term{reductive-group}{约化} $S$-群，$P$ 为 $G$ 的\Term{parabolic-subgroup}{抛物子群}。\Term{canonical}{典范}映射
\[\operatorname{rad}^{u}(P)(S)\longrightarrow\prod_i\operatorname{rad}^{u}(P)(\operatorname{Spec}\kappa(s_i))\]
为\Term{surjectivity}{满射}。
\end{SGAresult}
\begin{proof}
若 $S=\operatorname{Spec}(A)$，$\kappa(s_i)=A/\mathfrak p_i$，且 $\mathcal E$ 对应于\Term{projective-module}{投射}（从而\Term{flat}{平坦}）模 $E$，则须证映射
\[E\longrightarrow\prod_i E\otimes_A A/\mathfrak p_i\]
\Term{surjectivity}{满射}。只须在 $E=A$ 时证明，而此时是熟知的（参见 Bourbaki，\textit{Alg. Comm.}，第 II 章，§1，no.~2，命题 5）。
\end{proof}
\begin{SGAresult}{推论}{2.7}
设 $k$ 为无限域，$G$ 为\Term{reductive-group}{约化} $k$-群，$P$ 为 $G$ 的\Term{parabolic-subgroup}{抛物子群}；则 $\operatorname{rad}^{u}(P)(k)$ 在 $\operatorname{rad}^{u}(P)$ 中\Term{dense}{稠密}。
\end{SGAresult}
\begin{SGAresult}{推论}{2.8}
设 $S$ 为半\Term{local-scheme}{局部}\Term{scheme}{概形}，$\{s_i\}$ 为其\Term{closed-point}{闭点}集，$G$ 为\Term{reductive-group}{约化} $S$-群，$P$ 为 $G$ 的\Term{parabolic-subgroup}{抛物子群}，并对每个 $i$ 给定 $P_{s_i}$ 的\Term{levi-subgroup}{Levi 子群} $L_i$。存在 $P$ 的\Term{levi-subgroup}{Levi 子群} $L$，对每个 $i$ 诱导 $L_i$。
\end{SGAresult}
事实上，取 $P$ 的一个\Term{levi-subgroup}{Levi 子群} $L_0$（2.3）。对每个 $i$，取 $u_i\in\operatorname{rad}^{u}(P)(\operatorname{Spec}\kappa(s_i))$ 使 $\operatorname{int}(u_i)L_{0,s_i}=L_i$（1.8）；若 $u\in\operatorname{rad}^{u}(P)(S)$ 对每个 $i$ 诱导 $u_i$（2.6），则 $L=\operatorname{int}(u)L_0$ 即所求。
\begin{SGAresult}{推论}{2.9}
在 2.1 的情形中，再设 $H$ 为 $G$ 的\Term{subgroup}{子群}\Term{scheme}{概形}，在 $S$ 上\Term{smooth}{光滑}且\Term{finite-presentation}{有限表现}，\Term{fiber}{纤维}\Term{connected}{连通}，并使 $P\cap H$ 关于 $(\mathrm{fpqc})$ \Term{topology}{拓扑}\Term{local-scheme}{局部}包含 $G$ 的\Term{maximal-torus}{极大环面}。则对每个 $i\ge0$，存在 $\mathcal E_i$ 的一个\Term{local-scheme}{局部}为\Term{direct-summand}{直和因子}的\Term{submodule}{子模} $\mathcal F_i$，使\Term{isomorphism}{同构} $U_i/U_{i+1}\xrightarrow{\sim}\mathbf W(\mathcal E_i)$ 诱\Term{derived-group}{导群}\Term{isomorphism}{同构}
\[(U_i\cap H)/(U_{i+1}\cap H)\xrightarrow{\sim}\mathbf W(\mathcal F_i).\]
\end{SGAresult}
\SourcePage{289}
\begin{proof}
$H$ 是 $G$ 的 $(R)$ 型\Term{subgroup}{子群}（第 XXII 讲，5.2.1）。所证断言关于 $(\mathrm{fpqc})$ \Term{topology}{拓扑}\Term{local-scheme}{局部}，故可假设 $G$ 关于 $P\cap H$ 的一个\Term{maximal-torus}{极大环面}\Term{splitting}{分裂}；甚至可归约到 2.1.1 的情形，其中 $H$ 由 $R$ 的\Term{subset}{子集} $R'$ 定义。沿用该处记号，由第 XXII 讲，5.6.7(ii) 得 $U_i\cap H=\prod_{\alpha\in R_i\cap R'}U_\alpha$，故 $(U_i\cap H)/(U_{i+1}\cap H)$ 等同于 $\prod_{\alpha\in R',\,a(\alpha)=i+1}U_\alpha$，结论随之成立。
\end{proof}
\begin{SGAresult}{推论}{2.10}
在 2.9 的情形中，将 $\operatorname{rad}^{u}(P)$ 换成 $\operatorname{rad}^{u}(P)\cap H$，2.2、2.5、2.6、2.7 的结论仍成立。
\end{SGAresult}
\begin{SGAresult}{推论}{2.11}
\footnote{编者注：增添了推论 2.11，将在 4.5.1 及 6.17 中使用。}设 $G$ 为\Term{reductive-group}{约化} $S$-群，$P$ 为\Term{parabolic-subgroup}{抛物子群}，$H$ 为 $P$ 的 $(\mathrm{RC})$ 型\Term{subgroup}{子群}，且 $\operatorname{rad}^{u}(H)=\operatorname{rad}^{u}(P)\cap H$。则将 $P$ 换成 $H$，2.2 至 2.8 的陈述仍成立。
\end{SGAresult}
\SGAsection{3}{约化群的抛物子群概形}
\textbf{3.1.} 设 $E$ 为\Term{finite}{有限}\Term{twisted-constant}{扭常值} $S$-\Term{scheme}{概形}（第 X 讲，5.1）。考虑 $S$-\Term{functor}{函子} $\operatorname{Of}(E)$，其中 $\operatorname{Of}(E)(S')$ 是 $E_{S'}$ 的开闭子\Term{scheme}{概形}集（\Term{equivalence}{等价}地，$E_{S'}$ 的开\Term{closed-subset}{闭子集}集）。则 $\operatorname{Of}(E)$ 可由\Term{finite}{有限}\Term{twisted-constant}{扭常值} $S$-\Term{scheme}{概形}表示。事实上，若 $E=A_S$，$A$ 为\Term{finite}{有限}集，则立即有 $\operatorname{Of}(E)\simeq\mathcal P(A)_S$（$\mathcal P(A)$ 表示 $A$ 的幂集），再由开闭子\Term{scheme}{概形}的\Term{descent}{下降}得出结论。显然有
\[\operatorname{Of}(E_{S'})=\operatorname{Of}(E)_{S'},\qquad\operatorname{Of}(E\times_SE')=\operatorname{Of}(E)\times_S\operatorname{Of}(E').\]
\textbf{3.2.} 设 $S$ 为\Term{scheme}{概形}，$G$ 为\Term{reductive-group}{约化} $S$-群。$G$ 的\Term{parabolic-subgroup}{抛物子群}\Term{functor}{函子} $\operatorname{Par}(G)$ 定义为
\[\operatorname{Par}(G)(S')=\{G_{S'}\text{ 的抛物子群}\}.\]
特别地，$G\in\operatorname{Par}(G)(S)$，$\operatorname{Bor}(G)\subset\operatorname{Par}(G)$。我们要定义\Term{morphism}{态射}
\[t:\operatorname{Par}(G)\longrightarrow\operatorname{Of}(\operatorname{Dyn}(G))\]
使满足：
(i) $t$ 关于 $G$（对于\Term{isomorphism}{同构}）具有\Term{functorial}{函子性}，并与\Term{base-extension}{基扩张}交换。
(ii) 若 $(T,M,R,\Delta,\ldots)$ 为适合于\Term{parabolic-subgroup}{抛物子群} $P$ 的\Term{pinning}{钉扎}（1.11），则\Term{canonical-isomorphism}{典范同构} $\operatorname{Dyn}(G)\simeq\Delta_S$（第 XXIV 讲，3.4(iii)）将 $t(P)$ 变为 $\Delta(P)_S$（记号见 1.11、1.12）。
先取 $G$ 的\Term{parabolic-subgroup}{抛物子群} $P$ 及适合于 $P$ 的\Term{pinning}{钉扎} $(T,M,R,\Delta,\ldots)$，用 (ii) 定义 $t(P)$。这样构造的 $\operatorname{Dyn}(G)$ 的子\Term{scheme}{概形} $t(P)$ 不依赖于所选\Term{pinning}{钉扎}。事实上，若 $(T',M',R',\Delta',\ldots)$ 是另一个适合于 $P$ 的\Term{pinning}{钉扎}，则将第一个变为第二个的唯一\Term{inner-automorphism}{内自同构}
\SourcePage{290}
来自 $P$（1.15）；故\Term{canonical-isomorphism}{典范同构} $\Delta\xrightarrow{\sim}\Delta'$ 将 $\Delta(P)$ 变为 $\Delta'(P)$，得所述结论。
不再假设 $(G,P)$ 必可\Term{pinning}{钉扎}时，由 1.14 及 $\operatorname{Dyn}(G)$ 的定义（第 XXIV 讲，3.3），立即可由\Term{descent}{下降}定义 $\operatorname{Dyn}(G)$ 的唯一开闭子\Term{scheme}{概形} $t(P)$，使对每个令 $(G,P)_{S'}$ 可\Term{pinning}{钉扎}的 $S'\to S$，均有 $t(P)_{S'}=t(P_{S'})$。
\begin{SGAresult}{定理}{3.3}
设 $S$ 为\Term{scheme}{概形}，$G$ 为\Term{reductive-group}{约化} $S$-群，
\[t:\operatorname{Par}(G)\longrightarrow\operatorname{Of}(\operatorname{Dyn}(G))\]
为以上定义的\Term{morphism}{态射}。
(i) $G$ 的两个\Term{parabolic-subgroup}{抛物子群} $P,P'$ 关于 $(\mathrm{fpqc})$ \Term{topology}{拓扑}\Term{local-scheme}{局部}\Term{conjugacy}{共轭}（参见 1.3），当且仅当 $t(P)=t(P')$。
(ii) $\operatorname{Par}(G)$ \Term{representable-functor}{可表示}，且\Term{morphism}{态射} $t$ \Term{smooth}{光滑}、\Term{projective}{射影}、\Term{geometric-fibre}{几何纤维}整。
\end{SGAresult}
\begin{proof}
由 3.2(i) 及 $G$ 的\Term{inner-automorphism}{内自同构}在 $\operatorname{Dyn}(G)$ 上\Term{trivial-action}{平凡作用}（第 XXIV 讲，3.4(iv)），当 $P,P'$ \Term{conjugacy}{共轭}时确有 $t(P)=t(P')$。反之，设 $t(P)=t(P')$。为证它们在 $G$ 中关于 $(\mathrm{fpqc})$ \Term{topology}{拓扑}\Term{local-scheme}{局部}\Term{conjugacy}{共轭}，可先假设 $(G,P),(G,P')$ 可\Term{pinning}{钉扎}（1.14）；利用\Term{pinning}{钉扎}在 $G$ 中的\Term{conjugacy}{共轭}性（第 XXIV 讲，1.5），可假设 $G$ 有一个同时适合于 $P,P'$ 的\Term{pinning}{钉扎} $(T,M,R,\Delta)$。此时 $t(P)=t(P')$ 推出 $\Delta(P)=\Delta(P')$，故 $P=P'$（参见 1.4(v)）。(i) 已证。
证明 (ii)，沿用第 XXII 讲，5.11.5 的记号。\footnote{编者注：回顾该处记号，$\mathcal H_c=\mathcal H_c(G)$ 表示 $G$ 的 $(\mathrm{RC})$ 型\Term{subgroup}{子群}\Term{functor}{函子}，$\mathcal C\ell_c=\mathcal C\ell_c(G)$ 表示这类\Term{subgroup}{子群}的“\Term{conjugacy}{共轭}类”\Term{functor}{函子}，$c\ell:\mathcal H_c\to\mathcal C\ell_c$ 为\Term{canonical-projection}{典范投影}。}
有\Term{canonical}{典范}\Term{morphism}{态射} $\operatorname{Par}(G)\to\mathcal H_c$；显然（例如归约到可\Term{pinning}{钉扎}的情形）它\Term{embedding}{嵌入}下列\Term{cartesian-square}{笛卡尔方块}，其中竖\Term{arrow}{箭头}为\Term{monomorphism}{单态射}：
\[\begin{tikzcd}
\operatorname{Par}(G)\arrow[r,"t"]\arrow[d,hook]&\operatorname{Of}(\operatorname{Dyn}(G))\arrow[d,hook]\\
\mathcal H_c\arrow[r,"c\ell"]&\mathcal C\ell_c.
\end{tikzcd}\]
由所引处，$\mathcal H_c$ \Term{representable-functor}{可表示}，且 $c\ell$ \Term{smooth}{光滑}、\Term{quasi-projective-morphism}{拟射影}、\Term{finite-presentation}{有限表现}、\Term{geometric-fibre}{几何纤维}整，故 $t$ 也一样。
尚须证明 $t$ \Term{proper}{固有}。现在这已是关于 $(\mathrm{fpqc})$ \Term{topology}{拓扑}\Term{local-scheme}{局部}的断言，可归约到带\Term{pinning}{钉扎}的情形 $G=(G,T,M,R,\Delta,\ldots)$。此时 $\operatorname{Dyn}(G)\simeq\Delta_S$，只须对每个 $\Delta_1\subset\Delta$ 证明 $S$-\Term{scheme}{概形} $t^{-1}((\Delta_1)_S)$ 在 $S$ 上\Term{proper}{固有}。若 $P_1$ 为 $G$ 中包含 $T$ 且 $\Delta(P_1)=\Delta_1$ 的\Term{parabolic-subgroup}{抛物子群}，则由 (i)，\Term{set}{集合}上为 $g\mapsto\operatorname{int}(g)P_1$ 的\Term{morphism}{态射} $G\to\operatorname{Par}(G)$
\SourcePage{291}
诱导 $G/\operatorname{Norm}_G(P_1)$ 到 $t^{-1}((\Delta_1)_S)$ 的\Term{isomorphism}{同构}。由 1.2，$G/\operatorname{Norm}_G(P_1)=G/P_1$ 在 $S$ 上\Term{projective}{射影}。
\end{proof}
\begin{SGAresult}{定义}{3.4}
$\operatorname{Of}(\operatorname{Dyn}(G))$ 称为 $G$ 的\Term{parabolic-subgroup}{抛物子群}类型\Term{scheme}{概形}；$t(P)$ 称为 $P$ 的类型。
\end{SGAresult}
\begin{SGAresult}{推论}{3.5}
$S$-\Term{functor}{函子} $\operatorname{Par}(G)$ 可由在 $S$ 上\Term{smooth}{光滑}且\Term{projective}{射影}的 $S$-\Term{scheme}{概形}表示。\Term{resolution}{分解}
\[\operatorname{Par}(G)\longrightarrow\operatorname{Of}(\operatorname{Dyn}(G))\longrightarrow S\]
是\Term{structural-morphism}{结构态射} $\operatorname{Par}(G)\to S$ 的 \Term{stein-factorization}{Stein 分解}（EGA III，4.3.3）。
\end{SGAresult}
\begin{SGAresult}{推论}{3.6}
对每个 $t\in\operatorname{Of}(\operatorname{Dyn}(G))(S)$，$G$ 的 $t$ 型\Term{parabolic-subgroup}{抛物子群} $S$-\Term{scheme}{概形}
\[\operatorname{Par}_t(G)=t^{-1}(t)\]
在 $S$ 上\Term{smooth}{光滑}且\Term{projective}{射影}，是 $G$ 的\Term{homogeneous-space}{齐性空间}。若 $P$ 为 $G$ 的\Term{parabolic-subgroup}{抛物子群}，则有\Term{canonical-isomorphism}{典范同构} $G/P\xrightarrow{\sim}\operatorname{Par}_{t(P)}(G)$。有 $\operatorname{Par}_\varnothing(G)=\operatorname{Bor}(G)$，$\operatorname{Par}_{\operatorname{Dyn}(G)}(G)\xrightarrow{\sim}S$。
\end{SGAresult}
\textbf{注记 3.7。} $S$-\Term{scheme}{概形} $\operatorname{Of}(\operatorname{Dyn}(G))$ 带有自然序结构（子\Term{scheme}{概形}间的\Term{dominant}{支配}关系，此处即\Term{set}{集合}包含关系）。这个序结构为格；特别地，$\operatorname{Dyn}(G)_{S'}$ 的两个开闭子\Term{scheme}{概形}的\Term{infimum}{下确界}显然是它们的交。还应指出，若 $B$ 为 $G$ 的\Term{borel-subgroup}{Borel 子群}，可定义 $G$ 中包含 $B$ 的\Term{parabolic-subgroup}{抛物子群}\Term{functor}{函子} $X$。由断言 $P\subset Q\Rightarrow t(P)\subset t(Q)$ 及下述引理，$t$ 诱导的\Term{morphism}{态射} $X\to\operatorname{Of}(\operatorname{Dyn}(G))$ 是\Term{isomorphism}{同构}（对“有序\Term{scheme}{概形}”结构而言）。
\begin{SGAresult}{引理}{3.8}
设 $S$ 为\Term{scheme}{概形}，$G$ 为\Term{reductive-group}{约化} $S$-群，$P$ 为 $G$ 的\Term{parabolic-subgroup}{抛物子群}，$t'$ 为 $\operatorname{Of}(\operatorname{Dyn}(G))$ 在 $S$ 上的\Term{section}{截面}，且 $t(P)\subset t'$。存在唯一的 $G$ 的\Term{parabolic-subgroup}{抛物子群} $P'$，使 $P'\supset P$ 且 $t(P')=t'$。
\end{SGAresult}
\begin{proof}
\Term{group-extension}{扩张}\Term{base}{基底}后，可假设 $P$ 包含 $G$ 的\Term{borel-subgroup}{Borel 子群} $B$。此时唯一性由 1.17 推出。存在性可在\Term{splitting}{分裂}情形中证明，而此时显然，参见第 1 节。
\end{proof}
\textbf{注记 3.8.1。} (i) 将 3.8 的包含关系倒转所得的类似断言显然错误。例如它会推出每个 $A_1$ 型群都有\Term{borel-subgroup}{Borel 子群}，而事实并非如此，参见第 XX 讲，第 5 节。
(ii) 以上立即表明，$t(P)\subset t(Q)$ 意味着关于 $(\mathrm{fpqc})$ 或 $(\text{ét})$ \Term{local-scheme}{局部}地，$P$ 与 $Q$ 的一个\Term{subgroup}{子群}\Term{conjugacy}{共轭}（且只须在\Term{geometric-fibre}{几何纤维}上检验）。第 5 节将说明可将\Term{etale}{平展}\Term{topology}{拓扑}换成 Zariski \Term{topology}{拓扑}。
\SourcePage{292}
\textbf{3.9.} 以上讨论可对\Term{critical-subgroup}{临界子群}重复进行。回顾（第 XXII 讲，5.10.4、5.10.5），\Term{reductive-group}{约化群} $G$ 的\Term{reductive-group}{约化子群} $H$ 称为\Term{critical-subgroup}{临界子群}，若 $H=\operatorname{Centr}_G(\operatorname{rad}(H))$；$G$ 的子\Term{torus}{环面} $Q$ 称为 $C$-\Term{critical-torus}{临界环面}，若 $Q=\operatorname{rad}(\operatorname{Centr}_G(Q))$。\Term{critical-subgroup}{临界子群}与 $C$-\Term{critical-torus}{临界环面}\footnote{编者注：此处将“\Term{critical-torus}{临界环面}”改为“$C$-\Term{critical-torus}{临界环面}”，参见所引处。以下将简称“\Term{critical-torus}{临界环面}”。}一一对应（由 $H\mapsto\operatorname{rad}(H)$ 及 $Q\mapsto\operatorname{Centr}_G(Q)$）。
若 $(G,T,M,R)$ 为\Term{splitting}{分裂} $S$-群，则与 $R$ 的\Term{subset}{子集} $R'$ 对应的 $G$ 中包含 $T$ 的\Term{subgroup}{子群}（第 XXII 讲，5.4.2）为\Term{critical-subgroup}{临界子群}，当且仅当 $R'$ 为“向量性的”（即 $R$ 与 $M\otimes\mathbf Q$ 中某向量子\Term{space}{空间}的交），参见第 XXII 讲，5.10.6。
若 $G$ 为任意\Term{reductive-group}{约化} $S$-群，可如第 XXII 讲，5.11.5，定义\Term{finite}{有限}\Term{etale}{平展} $S$-\Term{scheme}{概形} $\mathcal C\ell_{\mathrm{crit}}$；在\Term{splitting}{分裂}情形，它是 $R$ 的向量性\Term{subset}{子集}集模 \Term{weyl-group}{Weyl 群}\Term{action}{作用}所得\Term{set}{集合}对应的常值\Term{scheme}{概形}。若 $\operatorname{Crit}(G)$ 为 $G$ 的“\Term{critical-subgroup}{临界子群}\Term{functor}{函子}”，则有\Term{canonical}{典范}\Term{morphism}{态射} $c\ell:\operatorname{Crit}(G)\to\mathcal C\ell_{\mathrm{crit}}$，位于\Term{cartesian-diagram}{笛卡尔图}中：\footnote{编者注：$\mathcal H_c$ 及 $\mathcal C\ell_c$ 的记号参见 3.3 的脚注 (9)。}
\[\begin{tikzcd}
\operatorname{Crit}(G)\arrow[r,"c\ell"]\arrow[d,hook]&\mathcal C\ell_{\mathrm{crit}}\arrow[d,hook]\\
\mathcal H_c\arrow[r,"c\ell"]&\mathcal C\ell_c.
\end{tikzcd}\]
\begin{SGAresult}{命题}{3.10}
设 $S$ 为\Term{scheme}{概形}，$G$ 为\Term{reductive-group}{约化} $S$-群，$\operatorname{Crit}(G)$ 为其\Term{critical-subgroup}{临界子群}\Term{functor}{函子}，$\mathcal C\ell_{\mathrm{crit}}$ 及 $c\ell:\operatorname{Crit}(G)\to\mathcal C\ell_{\mathrm{crit}}$ 为上述\Term{finite}{有限}\Term{etale}{平展} $S$-\Term{scheme}{概形}及\Term{morphism}{态射}。
(i) $G$ 的\Term{critical-subgroup}{临界子群} $H,H'$ 关于 $(\mathrm{fpqc})$ \Term{topology}{拓扑}\Term{local-scheme}{局部}\Term{conjugacy}{共轭}，当且仅当 $c\ell(H)=c\ell(H')$。
(ii) $\operatorname{Crit}(G)$ \Term{representable-functor}{可表示}，且 $c\ell$ \Term{smooth}{光滑}、\Term{affine}{仿射}，\Term{geometric-fibre}{几何纤维}整。
\end{SGAresult}
\begin{proof}
除“$c\ell$ \Term{affine}{仿射}”外，证明同 3.3。只须证明 $\operatorname{Crit}(G)$ 在 $S$ 上\Term{affine}{仿射}。$\operatorname{Crit}(G)$ 自然等同于 $G$ 的\Term{critical-torus}{临界环面} $S$-\Term{functor}{函子}，故有\Term{canonical}{典范}\Term{monomorphism}{单态射}
\[\operatorname{Crit}(G)\longrightarrow\mathcal M,\]
其中 $\mathcal M$ 为 $G$ 的\Term{multiplicative-type}{乘法型}\Term{subgroup}{子群}\Term{scheme}{概形}（第 XI 讲，4.1）。由第 XII 讲，5.3，为证 $\operatorname{Crit}(G)$ 在 $S$ 上\Term{affine}{仿射}，只须说明此\Term{morphism}{态射}为开\Term{closed-immersion}{闭浸入}；或者作\Term{base-change}{基变换} $\mathcal M\to S$ 后，只须证：若 $Q$ 为\Term{reductive-group}{约化群} $G$ 的\Term{multiplicative-type}{乘法型}\Term{subgroup}{子群}，则使 $Q_{S'}$ 为 $G_{S'}$ 的\Term{critical-torus}{临界环面}的那些 $S'\to S$，恰为经由 $S$ 的某一开闭子\Term{scheme}{概形}\Term{resolution}{分解}的\Term{morphism}{态射}。$Q$ 为\Term{critical-torus}{临界环面}，就是说：
\SourcePage{293}
(1) $Q$ 为\Term{torus}{环面}；
(2) 在 $Q$ 为\Term{torus}{环面}的条件下，同样为\Term{torus}{环面}的 $\operatorname{rad}(\operatorname{Centr}_G(Q))$ 与 $Q$ \Term{relative-dimension}{相对维数}相同。
而这两个条件确为所需的类型。
\end{proof}
\begin{SGAresult}{推论}{3.11}
$S$-\Term{functor}{函子} $\operatorname{Crit}(G)$ 可由在 $S$ 上\Term{smooth}{光滑}且\Term{affine}{仿射}的 $S$-\Term{scheme}{概形}表示。
\end{SGAresult}
\begin{SGAresult}{推论}{3.12}
设 $H$ 为\Term{reductive-group}{约化} $S$-群 $G$ 的\Term{critical-subgroup}{临界子群}。则 $G/\operatorname{Norm}_G(H)$ 及 $G/H$ 可由在 $S$ 上\Term{affine}{仿射}且\Term{smooth}{光滑}的 $S$-\Term{scheme}{概形}表示。
\end{SGAresult}
第一个断言由 3.11 推出，第二个由第一个及第 XXII 讲，5.10.2 推出。
\begin{SGAresult}{推论}{3.13}
设 $Q$ 为\Term{reductive-group}{约化} $S$-群 $G$ 的子\Term{torus}{环面}。则 $G/\operatorname{Centr}_G(Q)$ 可由在 $S$ 上\Term{smooth}{光滑}且\Term{affine}{仿射}的 $S$-\Term{scheme}{概形}表示。若 $Q$ 为 $G$ 的\Term{critical-torus}{临界子环面}，则 $G/\operatorname{Norm}_G(Q)$ 也一样。
\end{SGAresult}
事实上，$H=\operatorname{Centr}_G(Q)$ 为\Term{critical-subgroup}{临界子群}（第 XXII 讲，5.10.5）；若 $Q$ 为\Term{critical-torus}{临界环面}，则 $\operatorname{Norm}_G(H)=\operatorname{Norm}_G(Q)$（同处，5.10.8）。
\textbf{3.14.} 由 $G$ 的\Term{parabolic-subgroup}{抛物子群}的\Term{levi-subgroup}{Levi 子群}的\Term{conjugacy}{共轭}性，存在唯一\Term{morphism}{态射}
\[u:\operatorname{Of}(\operatorname{Dyn}(G))\longrightarrow\mathcal C\ell_{\mathrm{crit}},\]
使对 $G$ 的每个\Term{parabolic-subgroup}{抛物子群} $P$ 及 $P$ 的每个\Term{levi-subgroup}{Levi 子群} $L$，均有 $c\ell(L)=u(t(P))$，且在任意\Term{base-change}{基变换}后仍成立。
\textbf{3.15.} 设 $S$ 为\Term{scheme}{概形}，$G$ 为\Term{reductive-group}{约化} $S$-群。考虑以下 $S$-\Term{functor}{函子}：\footnote{编者注：将 LT 改为 CT。}
\[\begin{aligned}
\mathrm{PL}(S')&=\{P\supset L\mid P\text{ 为 }G_{S'}\text{ 的抛物子群，}L\text{ 为 }P\text{ 的 Levi 子群}\};\\
\mathrm{PT}(S')&=\{P\supset T\mid P\text{ 为 }G_{S'}\text{ 的抛物子群，}T\text{ 为 }P\text{ 的极大环面}\};\\
\mathrm{CT}(S')&=\{C\supset T\mid C\text{ 为 }G_{S'}\text{ 的临界子群，}T\text{ 为 }C\text{ 的极大环面}\};\\
\mathrm{PLT}(S')&=\{P\supset L\supset T\mid(P,T)\in\mathrm{PT}(S'),\ L\text{ 为 }P\text{ 的 Levi 子群}\}.
\end{aligned}\]
这些\Term{functor}{函子}与先前引入的 $\operatorname{Par}(G),\operatorname{Crit}(G),\operatorname{Tor}(G)$ 之间有明显的\Term{morphism}{态射}，且有一个截去一角的立方体形\Term{commutative-diagram}{交换图}（见下图）。
\SourcePage{294}
\begin{center}
\begin{tikzpicture}[x=1cm,y=1cm,>=stealth, every node/.style={font=\small}]
\node (CT) at (3,5) {$\mathrm{CT}$};\node (PLT) at (7,5) {$\mathrm{PLT}$};
\node (Crit) at (0,3.7) {$\operatorname{Crit}(G)$};\node (PL) at (5.5,3.7) {$\mathrm{PL}$};
\node (Cl) at (0,2.2) {$\mathcal C\ell_{\mathrm{crit}}$};\node (Tor) at (3,2.2) {$\operatorname{Tor}(G)$};\node (PT) at (7,2.2) {$\mathrm{PT}$};
\node (S) at (1.5,.9) {$S$};\node (Of) at (1.5,-.6) {$\operatorname{Of}(\operatorname{Dyn}(G))$};\node (Par) at (5.5,-.6) {$\operatorname{Par}(G)$};
\draw[->] (PLT)--node[above] {$g$（平展）}(CT);
\draw[->] (CT)--node[above,sloped] {$a$（仿射）}(Crit);
\draw[->] (PLT)--node[left] {$b$（仿射）}(PL);
\draw[->] (PL)--node[above] {$f$（平展）}(Crit);
\draw[->] (CT)--node[left] {$k$（平展）}(Tor);
\draw[->] (PLT)--node[right] {$e$（同构）}(PT);
\draw[->] (PT)--node[above] {$h$（平展）}(Tor);
\draw[->] (Crit)--node[left] {$c\ell$（仿射）}(Cl);
\draw[->] (Cl)--node[left] {$q$（平展）}(S);
\draw[->] (Tor)--node[right] {$r$（仿射）}(S);
\draw[->] (Of)--node[right] {$p$（平展）}(S);
\draw[->] (Of) to[bend left=14] node[left] {$u$（平展）}(Cl);
\draw[->] (PL)--node[left] {$d$（仿射）}(Par);
\draw[->] (PT)--node[right] {$c$（仿射）}(Par);
\draw[->] (Par)--node[above] {$t$（射影）}(Of);
\end{tikzpicture}

图 3.15.1
\end{center}
\begin{SGAresult}{定理}{3.16}
（参见图 3.15.1。）
(i) 图中所有\Term{morphism}{态射}均\Term{smooth}{光滑}、\Term{surjectivity}{满射}、\Term{finite-presentation}{有限表现}。
(ii) 图中除 $t$ 之外的所有\Term{morphism}{态射}均\Term{affine}{仿射}；\Term{morphism}{态射} $t$ \Term{projective}{射影}。
(iii) 图中所有\Term{morphism}{态射}或为\Term{finite}{有限}\Term{etale}{平展}，或有整的\Term{geometric-fibre}{几何纤维}：\Term{morphism}{态射} $f,g,h,k,p,q,u$ 为\Term{finite}{有限}\Term{etale}{平展}；\Term{morphism}{态射} $a,b,c,d,r,c\ell$ 的\Term{geometric-fibre}{几何纤维}整；\Term{morphism}{态射} $e$ 为\Term{isomorphism}{同构}。
(iv) 方块 $(a,b,f,g)$ 为\Term{cartesian-square}{笛卡尔方块}。
\end{SGAresult}
\begin{proof}
由 1.6(ii)，$e$ 显然是\Term{isomorphism}{同构}。另一方面，(iv) 显然。
\Term{morphism}{态射} $a$ \Term{smooth}{光滑}、\Term{affine}{仿射}、\Term{geometric-fibre}{几何纤维}整：事实上，作\Term{base-change}{基变换} $\operatorname{Crit}(G)\to S$ 后，只须检验\Term{morphism}{态射} $\operatorname{Tor}(L_0)\to S$ 有这些性质，其中 $L_0$ 是泛\Term{critical-subgroup}{临界子群}。按定义 $L_0$ \Term{reductive-group}{约化}，故归结于第 XXII 讲，5.8.3。由 (iv)，$b$ 有相同性质。
由 1.9，$d$ 也\Term{smooth}{光滑}、\Term{affine}{仿射}、\Term{geometric-fibre}{几何纤维}整；故 $c=dbe^{-1}$ 也一样。\Term{morphism}{态射} $r$ 有这些性质（第 XXII 讲，5.8.3），$c\ell$ 也一样（3.10）。
\SourcePage{295}
另一方面，已证 $p,q$ 为\Term{finite}{有限}\Term{etale}{平展}\Term{surjectivity}{满射}（3.1、3.9）。若证 $f,k$ 为\Term{finite}{有限}\Term{etale}{平展}\Term{surjectivity}{满射}，则 $g$（由 (iv)）及 $h$（因 $h=kge^{-1}$）也有这些性质。$t$ 的所述性质已在 3.3(ii) 证明，故只剩证明 $f$（或 $k$）是\Term{finite}{有限}\Term{etale}{平展}\Term{surjectivity}{满射}。我们证明 $k$ 的情形，$f$ 的证明类似。
只须说明：若 $T$ 为 $G$ 的\Term{maximal-torus}{极大环面}，则 $G$ 中包含 $T$ 的\Term{critical-subgroup}{临界子群}\Term{functor}{函子} $C$ 可由\Term{fiber}{纤维}\Term{nonempty}{非空}的\Term{finite}{有限}\Term{etale}{平展} $S$-\Term{scheme}{概形}表示。可假设 $G$ 关于 $T$ \Term{splitting}{分裂}；令 $E$ 为 $R$（此\Term{splitting}{分裂}的\Term{root-system}{根系}）的向量性\Term{subset}{子集}集。则 $C$ 由 $E_S$ 表示（3.9），证明完成。
\end{proof}
\begin{SGAresult}{推论}{3.17}
上图所有\Term{functor}{函子}均可由在 $S$ 上\Term{smooth}{光滑}的 $S$-\Term{scheme}{概形}表示，且除 $\operatorname{Par}(G)$ 外，均在 $S$ 上\Term{affine}{仿射}。
\end{SGAresult}
\textbf{注记 3.18。} (i) \Term{morphism}{态射} $f:\mathrm{PL}\to\operatorname{Crit}(G)$ 为\Term{etale}{平展}\Term{surjectivity}{满射}，这表明 $G$ 的\Term{subgroup}{子群}是\Term{critical-subgroup}{临界子群}，当且仅当它关于\Term{etale}{平展}\Term{topology}{拓扑}\Term{local-scheme}{局部}地为 $G$ 的某个\Term{parabolic-subgroup}{抛物子群}的\Term{levi-subgroup}{Levi 子群}。
但一般不能认为 $f(S):\mathrm{PL}(S)\to\operatorname{Crit}(G)(S)$ \Term{surjectivity}{满射}：$G$ 的\Term{critical-subgroup}{临界子群} $C$ 完全可能不在 $S$ 上来自 $G$ 的\Term{parabolic-subgroup}{抛物子群}；例如\Term{maximal-torus}{极大环面}未必包含在\Term{borel-subgroup}{Borel 群}中（如 $\mathrm{SL}_2$ 的非\Term{splitting}{分裂}\Term{form}{形式}，参见第 XX 讲，第 5 节）。
(ii) 同样，$u:\operatorname{Of}(\operatorname{Dyn}(G))\to\mathcal C\ell_{\mathrm{crit}}$ 也可能不是\Term{isomorphism}{同构}：不同类型的\Term{parabolic-subgroup}{抛物子群}可有同类型的\Term{levi-subgroup}{Levi 子群}。例如 $A_2$ 型群中，其\Term{levi-subgroup}{Levi 子群}的\Term{semisimple-rank}{半单秩}为 $1$ 的\Term{parabolic-subgroup}{抛物子群}有两种类型（对应图的两个顶点），但秩 $1$ 的\Term{critical-subgroup}{临界子群}仅有一种类型。
$A_3$ 型群中的类似例子表明，即使在\Term{algebraically-closed-field}{代数闭域}上，不\Term{isomorphism}{同构}的\Term{parabolic-subgroup}{抛物子群}也可有同类型的\Term{levi-subgroup}{Levi 子群}。\footnote{编者注：即 $\mathrm{GL}_4$ 的\Term{parabolic-subgroup}{抛物子群}
\[P=\left\{\begin{pmatrix} *&*&*&*\\ *&*&*&*\\0&0&*&*\\0&0&0&*\end{pmatrix}\right\},\quad P'=\left\{\begin{pmatrix} *&*&*&*\\0&*&*&*\\0&*&*&*\\0&0&0&*\end{pmatrix}\right\}\]
不\Term{isomorphism}{同构}（因 $\operatorname{rad}^{u}(P)\not\simeq\operatorname{rad}^{u}(P')$），但它们的\Term{levi-subgroup}{Levi 子群}
\[L=\left\{\begin{pmatrix} *&*&0&0\\ *&*&0&0\\0&0&*&0\\0&0&0&*\end{pmatrix}\right\},\quad L'=\left\{\begin{pmatrix} *&0&0&0\\0&*&*&0\\0&*&*&0\\0&0&0&*\end{pmatrix}\right\}\]
由元素 $\begin{pmatrix}0&0&1&0\\1&0&0&0\\0&1&0&0\\0&0&0&1\end{pmatrix}$ \Term{conjugacy}{共轭}。}
以对\Term{principal-bundle}{主丛}理论的一个应用结束本节。
\begin{SGAresult}{引理}{3.20}
\footnote{编者注：保留了原著编号：没有 3.19。}设 $S$ 为\Term{scheme}{概形}，$G$ 为\Term{reductive-group}{约化} $S$-群，$F$ 为以 $G$ 为结构群的\Term{torsor}{主齐性丛}，$G^F$ 为对应的\Term{twisted-form}{扭形式}。等同 $\operatorname{Dyn}(G)$ 与
\SourcePage{296}
$\operatorname{Dyn}(G^F)$（第 XXIV 讲，3.5）。设 $P$ 为 $G$ 的\Term{parabolic-subgroup}{抛物子群}。有\Term{canonical-isomorphism}{典范同构}
\[F/P\xrightarrow{\sim}\operatorname{Par}_{t(P)}(G^F).\]
特别地，$F$ 的结构群可\Term{reduced}{约化}为 $P$，当且仅当 $G^F$ 有 $t(P)$ 型\Term{parabolic-subgroup}{抛物子群}。
\end{SGAresult}
证明与第 XXIV 讲，4.2.1 完全相同。
\textbf{3.21.} 若 $S$ 为\Term{scheme}{概形}，$G$ 为\Term{reductive-group}{约化} $S$-群，$t\in\operatorname{Of}(\operatorname{Dyn}(G))(S)$，则以 $\mathrm H^1_t(S,G)$ 表示 $\mathrm H^1(S,G)$ 中那些以 $G$ 为结构群的\Term{principal-bundle}{主丛} $F$ 的类构成的\Term{subset}{子集}，其中相关群 $G^F$ 有 $t$ 型\Term{parabolic-subgroup}{抛物子群}。若 $G$ 本身有 $t$ 型\Term{parabolic-subgroup}{抛物子群} $P$，则 $\mathrm H^1_t(S,G)$ 正是 $\mathrm H^1(S,P)$ 在 $\mathrm H^1(S,G)$ 中的像，故此像不依赖于 $P$ 的选择。
\SGAsection{4}{两个抛物子群的相对位置}
\textbf{4.1. 一个预备结果}
\begin{SGAresult}{引理}{4.1.1}
设 $k$ 为域，$G$ 为\Term{reductive-group}{约化} $k$-群，$P,P'$ 为 $G$ 的两个\Term{parabolic-subgroup}{抛物子群}。
(i) 则 $P\cap P'$ \Term{smooth}{光滑}，与 $G$ 有相同\Term{reductive-rank}{约化秩}，并包含 $G$ 的\Term{maximal-torus}{极大环面} $T$。
(ii) $P\cap P'$ 关于 $T$ 的\Term{root-set}{根集} $R'$ 是 $G$ 关于 $T$ 的\Term{root-set}{根集} $R$ 的\Term{closed-subset}{闭子集}。\footnote{编者注：增添了 (ii)，将在 4.5.1 中使用。}
\end{SGAresult}
\begin{proof}
先设 $k$ \Term{algebraically-closed}{代数闭}。取 $P$（或 $P'$）的\Term{borel-subgroup}{Borel 子群} $B$（或 $B'$）。已知存在 $g\in G(k)$ 使 $\operatorname{int}(g)B=B'$。另一方面，若 $T_0$ 为 $B$ 的\Term{maximal-torus}{极大环面}，置 $N=\operatorname{Norm}_G(T_0)$，则由 Bruhat 定理（Bible，§13.4，定理 3 的推论 1），$G(k)=B(k)N(k)B(k)$。故有 $b,b_1\in B(k)$、$n\in N(k)$ 使 $g=bnb_1$，从而 $\operatorname{int}(b)\operatorname{int}(n)B=B'$。于是
\[P\cap P'\supset B\cap B'=\operatorname{int}(b)(B\cap\operatorname{int}(n)B')\supset\operatorname{int}(b)T_0.\]
再设 $k$ 任意。由以上结果，$P_{\bar k}\cap P'_{\bar k}$ 包含 $G_{\bar k}$ 的\Term{maximal-torus}{极大环面}；由第 XXII 讲，5.4.5，推出 $(P\cap P')_{\bar k}$ “沿\Term{unit-section}{单位截面}”\Term{smooth}{光滑}，因为在域上，故处处\Term{smooth}{光滑}（第 VI$_{\mathrm A}$ 讲，1.3.1），从而 $P\cap P'$ \Term{smooth}{光滑}。由第 VI$_{\mathrm A}$ 讲，2.3.1，$P\cap P'$ 的单位\Term{connected-component}{连通分支} $(P\cap P')^0$ 为\Term{open-subgroup}{开子群}，在 $S$ 上\Term{smooth}{光滑}。对它应用第 XIV 讲，1.1，即得 (i)。
最后，$P$（或 $P'$）关于 $T$ 的\Term{root-set}{根集} $R_P$（或 $R_{P'}$）闭，且 $R'=R_P\cap R_{P'}$，得 (ii)。
\end{proof}
\textbf{注记 4.1.2。} 可证明（[BT65]，4.5），$P\cap P'$ \Term{connected}{连通}；\footnote{编者注：考虑到 4.5.1 所增细节，此处改动了原著（原著说“我们不使用这个事实”）。}将在 4.5.1 中使用。
\SourcePage{297}
\textbf{注记 4.1.3。} 上述引理对任意\Term{scheme}{概形}并不成立。例如，设 $G$ 为\Term{algebraically-closed-field}{代数闭域} $k$ 上的\Term{reductive-group}{约化群}，$B$ 为 $G$ 的\Term{borel-subgroup}{Borel 群}。取 $G=X$ 为\Term{base}{基底}，考虑 $G_X$ 的\Term{borel-subgroup}{Borel 子群} $B_1,B_2$，其中 $B_1=B_X$，$B_2=\operatorname{int}(g_0)B_1$，$g_0$ 是 $G_X$ 的\Term{canonical}{典范}\Term{section}{截面}（\Term{diagonal-section}{对角截面}）。对每个 $g\in X(k)$，$B_1\cap B_2$ 在 $g$ 处的\Term{fiber}{纤维}正是 $B\cap\operatorname{int}(g)B$。若 $B\ne G$，此\Term{fiber}{纤维}\Term{dimension}{维数}随 $g$ 变化，故 $B_1\cap B_2$ 不可能在 $X$ 上\Term{smooth}{光滑}。
\textbf{4.2. \Term{transversal-position}{横截位置}}
\begin{SGAresult}{定理}{4.2.1}
设 $S$ 为\Term{scheme}{概形}，$G$ 为\Term{reductive-group}{约化} $S$-群，$P,Q$ 为 $G$ 的两个\Term{parabolic-subgroup}{抛物子群}。关于 $(P,Q)$ 的下列条件\Term{equivalence}{等价}：
(i) $\operatorname{Lie}(P/S)+\operatorname{Lie}(Q/S)=\operatorname{Lie}(G/S)$。
(ii) \Term{canonical}{典范}\Term{morphism}{态射} $P\times_SQ\to G$ \Term{smooth}{光滑}。
(ii$'$) \Term{canonical}{典范}\Term{morphism}{态射} $P\to G/Q$ \Term{smooth}{光滑}。
(iii) \Term{canonical}{典范}\Term{morphism}{态射} $P\times_SQ\to G$ 是\Term{open-morphism}{开态射}。
(iii$'$) \Term{canonical}{典范}\Term{morphism}{态射} $P\to G/Q$ 是\Term{open-morphism}{开态射}。
(iv) \Term{canonical}{典范}\Term{morphism}{态射} $P\times_SQ\to G$ 逐\Term{fiber}{纤维}\Term{dominant}{支配}。
(iv$'$) \Term{canonical}{典范}\Term{morphism}{态射} $P\to G/Q$ 逐\Term{fiber}{纤维}\Term{dominant}{支配}。
(v) 对每个 $s\in S$，“$P_s\cap Q_s$ 的\Term{dimension}{维数}最小”，即
\[\dim(P_s\cap Q_s)=\dim P_s+\dim Q_s-\dim G_s.\]
(vi) 存在\Term{etale}{平展}\Term{topology}{拓扑}的\Term{covering-family}{覆盖族} $\{S_i\to S\}$，且对每个 $i$，存在 $P_{S_i}$ 的\Term{borel-subgroup}{Borel 子群} $B_i$ 及 $Q_{S_i}$ 的\Term{borel-subgroup}{Borel 子群} $B'_i$，使 $B_i\cap B'_i$ 是 $G_{S_i}$ 的\Term{maximal-torus}{极大环面}。
(vii) 存在\Term{etale}{平展}\Term{topology}{拓扑}的\Term{covering-family}{覆盖族} $\{S_i\to S\}$，且对每个 $i$，存在 $G_{S_i}$ 的\Term{splitting}{分裂} $(T_i,\ldots,R_i)$ 及 $R_i$ 的\Term{positive-root-system}{正根系} $R_i^+$，使 $P_{S_i}$（或 $Q_{S_i}$）是 $G_{S_i}$ 中包含 $T_i$、由 $R_i$ 的包含 $R_i^+$（或 $-R_i^+$）的\Term{subset}{子集} $R_i^{(1)}$（或 $R_i^{(2)}$）定义的 $(R)$ 型\Term{subgroup}{子群}（定义见第 XXII 讲，5.4.2、5.2.1）。
\end{SGAresult}
\begin{proof}
按以下逻辑图证明：
\[\begin{tikzcd}[column sep=large]
\text{(i)}\arrow[d,Rightarrow]&\text{(vii)}\arrow[l,Rightarrow]&\text{(vi)}\arrow[l,Rightarrow]\\
\text{(ii$'$)}\arrow[d,Rightarrow]\arrow[r,Rightarrow]&\text{(iii$'$)}\arrow[r,Rightarrow]&\text{(iv$'$)}\arrow[r,Rightarrow]&\text{(v)}\arrow[ul,Rightarrow]\\
\text{(ii)}\arrow[r,Rightarrow]&\text{(iii)}\arrow[r,Rightarrow]&\text{(iv)}\arrow[ur,Rightarrow]
\end{tikzcd}\]
显然 (ii)$\Rightarrow$(iii)、(ii$'$)$\Rightarrow$(iii$'$)。若 (iii) 成立，则 $P\times_SQ\to G$ 的\Term{set}{集合}像
\SourcePage{298}
是包含\Term{unit-section}{单位截面}的\Term{open-subset}{开集}；$G$ 的\Term{fiber}{纤维}\Term{connected}{连通}，故此像在每条\Term{fiber}{纤维}上\Term{dense}{稠密}，得 (iv)。同理有 (iii$'$)$\Rightarrow$(iv$'$)。
由\Term{cartesian-diagram}{笛卡尔图}
\[\begin{tikzcd}P\times_SQ\arrow[r]\arrow[d,"\operatorname{pr}_1"']&G\arrow[d]\\P\arrow[r]&G/Q\end{tikzcd}\]
得 (ii$'$)$\Rightarrow$(ii)。另一方面，由\Term{dimension}{维数}理论（参见 EGA IV$_2$，5.6.6），(iv) 或 (iv$'$) 蕴含 (v)。确可假设 $S=\operatorname{Spec}(k)$，$k$ 为\Term{algebraically-closed-field}{代数闭域}；直接计算表明，(iv)（或 (iv$'$)）的\Term{morphism}{态射}在 $G(k)$（或 $(G/Q)(k)$）的点处的每个\Term{nonempty}{非空}\Term{fiber}{纤维}均\Term{isomorphism}{同构}于 $P\cap Q$。
由第 XXII 讲，5.5.1(iv)、5.9.2，有 (vi)$\Rightarrow$(vii)。
有 (vii)$\Rightarrow$(i)：检验 $\operatorname{Lie}(P)+\operatorname{Lie}(Q)=\operatorname{Lie}(G)$ 可关于 $(\mathrm{fpqc})$ \Term{local-scheme}{局部}进行。若 (vii) 成立，可假设 $G$ \Term{splitting}{分裂}，$P\supset B_{R^+}$，$Q\supset B_{-R^+}$（沿用通常记号）；此时已有
\[\operatorname{Lie}(B_{R^+})+\operatorname{Lie}(B_{-R^+})=\operatorname{Lie}(G).\]
证明 (i) 蕴含 (ii$'$)。令 $u:P\to G/Q$ 为\Term{canonical}{典范}\Term{morphism}{态射}。因 $P,G/Q$ 在 $S$ 上\Term{smooth}{光滑}，为证 $u$ \Term{smooth}{光滑}，只须在 $u$ 的\Term{geometric-fibre}{几何纤维}上证明，故可假设 $S$ 为\Term{algebraically-closed-field}{代数闭域}的谱。\Term{morphism}{态射} $u$ 与显然的 $P$ \Term{action}{作用}相容（$u(pp')=pu(p')$），且 $P$ \Term{connected}{连通}，故只须检验 $u$ 在 $e\in G(k)$ 处\Term{smooth}{光滑}（第 VI$_{\mathrm A}$ 讲，1.3.1），即（SGA 1，第 II 讲，4.7）$u$ 在 $e$ 处的切映射\Term{surjectivity}{满射}。但此切映射自然等同于\Term{canonical}{典范}映射 $\operatorname{Lie}(P)\to\operatorname{Lie}(G)/\operatorname{Lie}(Q)$，由 (i) 它\Term{surjectivity}{满射}。
仅剩证明 (v)$\Rightarrow$(vi)。先设 $S$ 为\Term{algebraically-closed-field}{代数闭域}的谱。由 4.1.1，存在包含于 $P,Q$ 的\Term{maximal-torus}{极大环面} $T$；令 $R$（或 $R_1,R_2$）为 $G$（或 $P,Q$）关于 $T$ 的\Term{root-set}{根集}。
例如由第 XXII 讲，5.4.4、5.4.5，有
\[\begin{aligned}
\dim G&=\dim T+\operatorname{Card}(R),&\dim P&=\dim T+\operatorname{Card}(R_1),\\
\dim Q&=\dim T+\operatorname{Card}(R_2),&\dim(P\cap Q)&=\dim T+\operatorname{Card}(R_1\cap R_2).
\end{aligned}\]
所以 (v) \Term{equivalence}{等价}于 $\operatorname{Card}(R_1\cap R_2)=\operatorname{Card}(R_1)+\operatorname{Card}(R_2)-\operatorname{Card}(R)$，即 $R_1\cup R_2=R$。由第 XXII 讲，5.9.2、5.4.5，为证 (vi)，只须证明 $R_1\cap(-R_2)$ 包含 $R$ 的\Term{positive-root-system}{正根系}。故归结于：
\end{proof}
\SourcePage{299}
\begin{SGAresult}{引理}{4.2.2}
设 $R$ 为“\Term{root-system}{根系}”（例如第 XXI 讲意义下一个\Term{root-datum}{根数据}的\Term{root-set}{根集}）。设 $R_1,R_2$ 为 $R$ 的两个\Term{closed-subset}{闭子集}，各包含一个\Term{positive-root-system}{正根系}。若 $R_1\cup R_2=R$，则 $R_1\cap(-R_2)$ 包含一个\Term{positive-root-system}{正根系}。
\end{SGAresult}
\begin{proof}
$R_3=R_1\cap(-R_2)$ 显然闭；由第 XXI 讲，3.3.6，只须证明 $R_3\cup(-R_3)=R$。已知 $R_1\cup(-R_1)=R=R_2\cup(-R_2)$，利用以下初等事实即得结论：若 $A,A',B,B'$ 为\Term{set}{集合} $E$ 的四个\Term{subset}{子集}，且 $A\cup A'=B\cup B'=A\cup B=A'\cup B'=E$，则 $(A\cap B')\cup(A'\cap B)=E$。
\end{proof}
\begin{proof}[定理 4.2.1 证明续]
这完成了\Term{base}{基底}为\Term{algebraically-closed-field}{代数闭域}的谱时 (v)$\Rightarrow$(vi) 的证明。回到一般情形并假设 (v)。设 $s\in S$。由以上讨论，可取 $P_{\bar s}$（或 $Q_{\bar s}$）的\Term{borel-subgroup}{Borel 子群} $\overline B$（或 $\overline B'$），使 $\overline B\cap\overline B'$ 为 $G_{\bar s}$ 的\Term{maximal-torus}{极大环面}。
因 $P$ 的\Term{borel-subgroup}{Borel 子群} $S$-\Term{scheme}{概形} $\operatorname{Bor}(P)\simeq\operatorname{Bor}(P/\operatorname{rad}^{u}(P))$ \Term{smooth}{光滑}，可应用“Hensel 引理”（参见第 XI 讲，1.10），关于\Term{etale}{平展}\Term{topology}{拓扑}\Term{local-scheme}{局部}地（即将 $S$ 换成\Term{covering}{覆盖} $s$ 的\Term{etale}{平展} $S'\to S$，并将 $s$ 换成 $S'$ 中其\Term{fiber}{纤维}上的一点），假设 $P$ 有\Term{projection}{投影}到 $\overline B$ 的\Term{borel-subgroup}{Borel 子群} $B$；同样可假设 $Q$ 有\Term{projection}{投影}到 $\overline B'$ 的\Term{borel-subgroup}{Borel 子群} $B'$。因 $\overline B\cap\overline B'$ 为 $G_{\bar s}$ 的\Term{maximal-torus}{极大环面}，$S$ 有包含 $s$ 的\Term{open-subset}{开集} $U$，使 $B_U\cap B'_U$ 为 $G_U$ 的\Term{maximal-torus}{极大环面}（第 XXII 讲，5.9.4），得 (vi)。
\end{proof}
\begin{SGAresult}{定义}{4.2.3}
满足定理 4.2.1 的\Term{equivalence}{等价}条件 (i) 至 (vii) 的对 $(P,Q)$，称为处于\Term{transversal-position}{横截位置}。也称 $P$ 相对于 $Q$ 处于\Term{transversal-position}{横截位置}，或稍滥用地称 $P,Q$ （相互）处于\Term{transversal-position}{横截位置}。
\end{SGAresult}
由 (vi)，对\Term{borel-subgroup}{Borel 子群}而言，此定义与第 XXII 讲，5.9.1 的定义相同。
\begin{SGAresult}{推论}{4.2.4}
设 $S$ 为\Term{scheme}{概形}，$G$ 为\Term{reductive-group}{约化} $S$-群，$P,Q$ 为 $G$ 的两个\Term{parabolic-subgroup}{抛物子群}。
(i) $(P,Q)$ 处于\Term{transversal-position}{横截位置}，当且仅当对 $S$ 的每点 $s$，$(P_s,Q_s)$ 处于\Term{transversal-position}{横截位置}；若 $S'\to S$ \Term{surjectivity}{满射}且 $(P_{S'},Q_{S'})$ 处于\Term{transversal-position}{横截位置}，则 $(P,Q)$ 处于\Term{transversal-position}{横截位置}。
(ii) 存在 $S$ 的开子\Term{scheme}{概形} $U$，具有如下性质：\Term{morphism}{态射} $S'\to S$ 经由 $U$ \Term{resolution}{分解}，当且仅当 $(P_{S'},Q_{S'})$ 处于\Term{transversal-position}{横截位置}。
(iii) 考虑\Term{subfunctor}{子函子}
\[\operatorname{Gen}(G)\subset\operatorname{Par}(G)\times_S\operatorname{Par}(G),\quad\operatorname{Gen}(/Q)\subset\operatorname{Par}(G),\quad\operatorname{Gen}(P/Q)\subset G,\]
\SourcePage{300}
定义如下：对 $S'\to S$，$\operatorname{Gen}(G)(S')$ 为 $G_{S'}$ 中处于\Term{transversal-position}{横截位置}的\Term{parabolic-subgroup}{抛物子群}对的\Term{set}{集合}；$\operatorname{Gen}(/Q)(S')$ 为 $G_{S'}$ 中相对于 $Q_{S'}$ 处于\Term{transversal-position}{横截位置}的\Term{parabolic-subgroup}{抛物子群}集；$\operatorname{Gen}(P/Q)(S')$ 为使 $\operatorname{int}(g)P_{S'}$ 相对于 $Q_{S'}$ 处于\Term{transversal-position}{横截位置}的 $g\in G(S')$ 所成的\Term{set}{集合}。
每个\Term{functor}{函子}均可由相应 $S$-\Term{scheme}{概形} $\operatorname{Par}(G)\times_S\operatorname{Par}(G)$、$\operatorname{Par}(G)$、$G$ 中在 $S$ 上\Term{universally-schematically-dense}{普遍概形稠密}的开子\Term{scheme}{概形}表示\footnote{编者注：将“相对\Term{dense}{稠密}”改为“在 $S$ 上\Term{universally-schematically-dense}{普遍概形稠密}”，参见 EGA IV$_3$，定义 11.10.8。}（参见第 XVIII 讲，§1）。
\end{SGAresult}
\begin{proof}
(i) 由 4.2.1(i) 对“\Term{transversal-position}{横截位置}”的刻画立即推出。为证 (ii)，取 $U=S-\operatorname{Supp}(\operatorname{Coker}u)$，其中 $u$ 是\Term{canonical}{典范}\Term{morphism}{态射} $\operatorname{Lie}(P)\oplus\operatorname{Lie}(G)\to\operatorname{Lie}(G)$。
因有\Term{cartesian-diagram}{笛卡尔图}
\[\begin{tikzcd}
G\arrow[r,"f"]&\operatorname{Par}(G)&\operatorname{Par}(G)\arrow[r,"f'"]&\operatorname{Par}(G)\times_S\operatorname{Par}(G)\\
\operatorname{Gen}(P/Q)\arrow[u,hook]\arrow[r]&\operatorname{Gen}(/Q)\arrow[u,hook]&\operatorname{Gen}(/Q)\arrow[u,hook]\arrow[r]&\operatorname{Gen}(G)\arrow[u,hook]
\end{tikzcd}\]
（$f(g)=\operatorname{int}(g)P$，$f'(R)=(R,Q)$），只须对 $\operatorname{Gen}(G)$ 检验 (iii)。令 $P_0$ 为 $G_{\operatorname{Par}(G)}$ 的\Term{canonical}{典范}\Term{parabolic-subgroup}{抛物子群}，置
\[X=\operatorname{Par}(G)\times_S\operatorname{Par}(G),\qquad P=\operatorname{pr}_1^*(P_0),\quad Q=\operatorname{pr}_2^*(P_0).\]
对 $G_X$ 的\Term{parabolic-subgroup}{抛物子群} $P,Q$ 应用 (ii)，得到 $X$ 的开子\Term{scheme}{概形} $U$，直接检验可知它等同于 $\operatorname{Gen}(G)$。尚须检验\Term{dense}{稠密}性，这可在\Term{geometric-fibre}{几何纤维}上进行\footnote{编者注：参见 EGA IV$_3$，11.10.10。}，故可假设 $S=\operatorname{Spec}(k)$，$k$ \Term{algebraically-closed}{代数闭}。因 $\operatorname{Par}(G)$ \Term{smooth}{光滑}，只须说明 $\operatorname{Gen}(G)$ 与 $\operatorname{Par}(G)\times_S\operatorname{Par}(G)$ 的每个\Term{irreducible-component}{不可约分支}相交；换言之，由 3.3，只须说明对 $t,t'\in\operatorname{Of}(\operatorname{Dyn}(G))(S)$，存在处于\Term{transversal-position}{横截位置}的 $(P,P')$，使 $t(P)=t,t(P')=t'$。这直接成立：选取 $G$ 的\Term{borel-subgroup}{Borel 子群}对 $(B,B')$，使 $B\cap B'$ 为\Term{maximal-torus}{极大环面}（\Term{splitting}{分裂} $G$，再应用第 XXII 讲，5.9.2），然后应用 3.8 构造 $P\supset B,P'\supset B'$，满足 $t(P)=t,t(P')=t'$；于是 $P,P'$ 有所需类型，且由 4.2.1(vi) 处于\Term{transversal-position}{横截位置}。
\end{proof}
\begin{SGAresult}{推论}{4.2.5}
设 $S$ 为\Term{scheme}{概形}，$G$ 为\Term{reductive-group}{约化} $S$-群，$P,Q$ 为 $G$ 的两个\Term{parabolic-subgroup}{抛物子群}，$(P,Q)$ 处于\Term{transversal-position}{横截位置}。
(i) 设 $P',Q'$ 为 $G$ 的两个\Term{parabolic-subgroup}{抛物子群}，分别与 $P,Q$ 类型相同。$(P',Q')$ 处于\Term{transversal-position}{横截位置}，当且仅当它关于\Term{etale}{平展}\Term{topology}{拓扑}\Term{local-scheme}{局部}与 $(P,Q)$ \Term{conjugacy}{共轭}。（注意：第 5 节将说明可将\Term{etale}{平展}\Term{topology}{拓扑}换成 Zariski \Term{topology}{拓扑}。）
\SourcePage{301}
(ii) \Term{canonical}{典范}\Term{morphism}{态射} $P\times_SQ\to G$ 诱导\Term{smooth}{光滑}\Term{surjectivity}{满射} $P\times_SQ\to\operatorname{Gen}(Q/P)$ 及\Term{isomorphism}{同构}
\[(P\times_SQ)/(P\cap Q)\simeq\operatorname{Gen}(Q/P),\]
其中 $P\cap Q=R$ 在 $P\times_SQ$ 上按 $(p,q)r=(pr,r^{-1}q)$ \Term{action}{作用}。
(iii) \Term{canonical}{典范}\Term{morphism}{态射} $P\to\operatorname{Par}_{t(Q)}(G)$（\Term{set}{集合}上由 $p\mapsto\operatorname{int}(p)Q$ 定义）诱导\Term{smooth}{光滑}\Term{surjectivity}{满射} $P\to\operatorname{Gen}(/P)\cap\operatorname{Par}_{t(Q)}(G)$ 及\Term{isomorphism}{同构}
\[P/(P\cap Q)\simeq\operatorname{Gen}(/P)\cap\operatorname{Par}_{t(Q)}(G).\]
(iv) \Term{canonical}{典范}\Term{morphism}{态射} $G\to\operatorname{Par}_{t(P)}(G)\times_S\operatorname{Par}_{t(Q)}(G)$（\Term{set}{集合}上由 $g\mapsto(\operatorname{int}(g)P,\operatorname{int}(g)Q)$ 定义）诱导\Term{smooth}{光滑}\Term{surjectivity}{满射}
\[G\longrightarrow\operatorname{Gen}(G)\cap\bigl(\operatorname{Par}_{t(P)}(G)\times_S\operatorname{Par}_{t(Q)}(G)\bigr)\]
及\Term{isomorphism}{同构}
\[G/(P\cap Q)\simeq\operatorname{Gen}(G)\cap\bigl(\operatorname{Par}_{t(P)}(G)\times_S\operatorname{Par}_{t(Q)}(G)\bigr).\]
\end{SGAresult}
\begin{proof}
证明 (i)。充分性显然；证明必要性。设 $(P',Q')$ 处于\Term{transversal-position}{横截位置}。因 $P,P'$ 关于\Term{etale}{平展}\Term{topology}{拓扑}\Term{local-scheme}{局部}\Term{conjugacy}{共轭}，可假设 $P=P'$。故只须证明：若 $G$ 的两个\Term{parabolic-subgroup}{抛物子群} $Q,Q'$ 均相对于 $P$ 处于\Term{transversal-position}{横截位置}且类型相同，则它们关于\Term{etale}{平展}\Term{topology}{拓扑}\Term{local-scheme}{局部}由 $P$ 的\Term{section}{截面}\Term{conjugacy}{共轭}。利用 4.2.1(vi)，可假设 $P,P,Q,Q'$ 分别有\Term{borel-subgroup}{Borel 子群} $B,B',B_1,B'_1$，使 $B\cap B_1=T$、$B'\cap B'_1=T'$ 是 $G$ 的\Term{maximal-torus}{极大环面}。但 $P$ 的 Killing 对 $(B,T),(B',T')$ 关于\Term{etale}{平展}\Term{topology}{拓扑}在 $P$ 中\Term{local-scheme}{局部}\Term{conjugacy}{共轭}（1.16），故可假设 $B=B',T=T'$；此时由第 XXII 讲，5.9.2 得 $B_1=B'_1$，再由 3.8 得 $Q=Q'$。
(ii)、(iii)、(iv) 的证明平行。以 (ii) 为例：设 $g\in\operatorname{Gen}(Q/P)(S)$，即 $g\in G(S)$ 使 $\operatorname{int}(g)Q$ 相对于 $P$ 处于\Term{transversal-position}{横截位置}。由以上证明，$Q,\operatorname{int}(g)Q$ 关于\Term{etale}{平展}\Term{topology}{拓扑}\Term{local-scheme}{局部}由 $P$ 的\Term{section}{截面}\Term{conjugacy}{共轭}。\Term{local-scheme}{局部}可假设有 $p\in P(S)$ 使 $\operatorname{int}(g)Q=\operatorname{int}(p)Q$，故 $p^{-1}g\in\operatorname{Norm}_G(Q)(S)=Q(S)$；于是存在 $q\in Q(S)$ 使 $g=pq$。这证明 (ii) 的\Term{morphism}{态射}是\Term{etale}{平展}\Term{topology}{拓扑}的\Term{covering}{覆盖}。结合 4.2.1(ii)，它为\Term{smooth}{光滑}\Term{surjectivity}{满射}。另一方面，直接论证表明，\Term{morphism}{态射} $P\times_SQ\to G$ 在 $P\times_SQ$ 上定义的\Term{equivalence-relation}{等价关系}，正是群 $R=P\cap Q$ 按 $(p,q)r=(pr,r^{-1}q)$ \Term{action}{作用}所对应的\Term{equivalence-relation}{等价关系}，得 (ii) 最后一项（例如因为\Term{smooth}{光滑}\Term{surjectivity}{满射}为\Term{effective-epimorphism}{有效满态射}）。
\end{proof}
\textbf{注记 4.2.6。} 若 $P,Q$ 处于\Term{transversal-position}{横截位置}，常以 $P\cdot Q$ 表示 $G$ 的\Term{open-subset}{开集} $\operatorname{Gen}(Q/P)$；4.2.5(ii) 说明此记号合理。
\begin{SGAresult}{命题}{4.2.7}
设 $S$ 为\Term{scheme}{概形}，$G$ 为\Term{reductive-group}{约化} $S$-群，$P,Q$ 为 $G$ 的两个\Term{parabolic-subgroup}{抛物子群}，$(P,Q)$ 处于\Term{transversal-position}{横截位置}。
\SourcePage{302}
(i) 群 $P\cap Q$ 在 $S$ 上\Term{smooth}{光滑}（由 4.1.2，其\Term{fiber}{纤维}实际\Term{connected}{连通}）。引入 $(P\cap Q)^0$（参见第 VI$_{\mathrm B}$ 讲，3.10）；它是 $G$ 的 $(\mathrm{RC})$ 型\Term{subgroup}{子群}（第 XXII 讲，5.11.1），其\Term{unipotent-radical}{幂幺根基}（同处，5.11.4）\Term{resolution}{分解}为\Term{direct-product}{直积}
\[\operatorname{rad}^{u}((P\cap Q)^0)=(\operatorname{rad}^{u}(P)\cap Q)\times_S(P\cap\operatorname{rad}^{u}(Q)).\]
(ii) 若 $S$ \Term{affine}{仿射}，则 $\mathrm H^1(S,\operatorname{rad}^{u}(P\cap Q)^0)=0$。若 $S$ 半\Term{local-scheme}{局部}，则 $P\cap Q$ 包含 $G$ 的\Term{maximal-torus}{极大环面}。
\end{SGAresult}
\begin{proof}
由 4.2.1(ii) 及\Term{cartesian-diagram}{笛卡尔图}
\[\begin{tikzcd}P\times_SQ\arrow[r]&G\\P\cap Q\arrow[u]\arrow[r]&S\arrow[u]\end{tikzcd}\]
知 $P\cap Q$ \Term{smooth}{光滑}。为检验关于 $(P\cap Q)^0$ 的所述断言，可关于\Term{etale}{平展}\Term{topology}{拓扑}\Term{local-scheme}{局部}论证；由 4.2.1(vii)，可假设已选择 $G$ 的\Term{splitting}{分裂} $(G,T,M,R)$，使 $P,Q$ 包含 $T$，分别由 $R$ 的\Term{subset}{子集} $R_1,R_2$ 定义，$R_1$ 包含\Term{positive-root-system}{正根系} $R^+$，$R_2$ 包含相反\Term{root-system}{根系} $R^-$。令 $\Delta$ 为 $R^+$ 的\Term{simple-root-system}{单根集}，记
\[A_1=\Delta\cap(-R_1),\quad A_2=\Delta\cap R_2,\quad A=A_1\cap A_2.\]
由 1.4(v)、1.12，
\[\begin{aligned}
R_1&=R^+\cup(R^-\cap(-\mathbf N A_1)),&R_2&=(R^+\cap\mathbf N A_2)\cup R^-,\\
\operatorname{rad}^{u}(P)&=\prod_{\alpha\in R_1,\,\alpha\notin-R_1}U_\alpha,&\operatorname{rad}^{u}(Q)&=\prod_{\alpha\in R_2,\,\alpha\notin-R_2}U_\alpha.
\end{aligned}\]
由第 XXII 讲，5.6.7，
\[\begin{aligned}
\operatorname{rad}^{u}(P)\cap Q&=\prod_{\alpha\in R_1\cap R_2,\,\alpha\notin-R_1}U_\alpha=\prod_{\alpha\in K_2}U_\alpha,\\
\operatorname{rad}^{u}(Q)\cap P&=\prod_{\alpha\in R_1\cap R_2,\,\alpha\notin-R_2}U_\alpha=\prod_{\alpha\in K_1}U_\alpha,
\end{aligned}\]
其中 $K_2$ 为\Term{representable}{可表}成 $A_2$ 中元素的线性组合、但不能表成 $A$ 中元素的线性组合的\Term{positive-root}{正根}集；$K_1$ 为\Term{representable}{可表}成 $A_1$ 中元素的线性组合、但不能表成 $A$ 中元素的线性组合的\Term{negative-root}{负根}集。显然若 $\alpha\in K_2,\beta\in K_1$，则 $\alpha+\beta$ 既不是根也不为零，故上述两个群交换。
另一方面，由第 XXII 讲，5.4.5，$H=(P\cap Q)^0$ 由\Term{root-set}{根集} $R_1\cap R_2$ 定义，其中
\[R_1\cap R_2=(R^+\cap\mathbf N A_2)\cup(R^-\cap(-\mathbf N A_1)).\]
\SourcePage{303}
因 $R_1\cap R_2$ 闭，按定义（第 XXII 讲，5.11.1）$H$ 为 $(\mathrm{RC})$ 型，且由该处 5.11.3、5.11.4，
\[\operatorname{rad}^{u}(H)=\prod_{\alpha\in K}U_\alpha,\]
其中 $K=\{\alpha\in R_1\cap R_2\mid\alpha\notin-(R_1\cap R_2)\}$。$R_1\cap R_2$ 的对称部分显然为 $R\cap\mathbf ZA$，故立即可见 $K=K_1\cup K_2$，(i) 证毕。
(ii) 第一项由 (i) 及 2.10 推出。证明第二项。$(P\cap Q)^0/\operatorname{rad}^{u}((P\cap Q)^0)$ \Term{reductive-group}{约化}，故当\Term{base}{基底}半\Term{local-scheme}{局部}时有\Term{maximal-torus}{极大环面} $\overline T$。$\overline T$ 在 $(P\cap Q)^0$ 中的\Term{inverse-image}{逆像}是 $G$ 的一个\Term{fiber}{纤维}\Term{solvable}{可解}的 $(R)$ 型\Term{subgroup}{子群} $N$，且 $N^u=\operatorname{rad}^{u}((P\cap Q)^0)$（第 XXII 讲，5.6.9）。$N$ 的\Term{maximal-torus}{极大环面}\Term{scheme}{概形}是以 $N^u$ 为结构群的\Term{torsor}{主齐性丛}（同处，5.6.13），因 $\mathrm H^1(S,N^u)=0$，故有\Term{section}{截面}。
\end{proof}
\textbf{4.3. \Term{opposite-parabolic}{相反抛物子群}}
\textbf{4.3.1.} 若 $G$ 为\Term{reductive-group}{约化} $S$-群，第 XXIV 讲，3.16.6 定义了 $G$ 的一个阶 $\le2$ 的\Term{canonical}{典范}“\Term{outer-automorphism}{外自同构}”，\footnote{编者注：将“阶为 $2$”改成“阶 $\le2$”，因为 $s_G$ 可平凡（例如 $G$ 为 $A_1,B_n,C_n,\ldots$ 型时）。}从而定义 $\operatorname{Dyn}(G)$ 的阶 $\le2$ 的\Term{canonical}{典范}\Term{automorphism}{自同构} $s_G$，因而也定义 $\operatorname{Of}(\operatorname{Dyn}(G))$ 的阶 $\le2$ 的\Term{automorphism}{自同构}，仍记为 $s_G$ 或简称 $s$。两个\Term{parabolic-subgroup}{抛物子群}类型 $t,t'\in\operatorname{Of}(\operatorname{Dyn}(G))(S)$ 称为相反的，若 $t=s_G(t')$。
\begin{SGAresult}{定理}{4.3.2}
设 $S$ 为\Term{scheme}{概形}，$G$ 为\Term{reductive-group}{约化} $S$-群，$P$ 为 $G$ 的\Term{parabolic-subgroup}{抛物子群}。
(a) 若 $L$ 为 $P$ 的\Term{levi-subgroup}{Levi 子群}，则存在唯一的 $G$ 的\Term{parabolic-subgroup}{抛物子群} $P'$，使 $P\cap P'=L$。
(b) 对 $G$ 的每个\Term{parabolic-subgroup}{抛物子群} $Q$，下列条件\Term{equivalence}{等价}：
(i) 对每个 $s\in S$，$((P\cap Q)_s)^0$（由 4.1.1 它\Term{smooth}{光滑}）为\Term{reductive-group}{约化群}。
(ii) $P\cap Q$ 是 $P,Q$ 的共同\Term{levi-subgroup}{Levi 子群}。
(iii) $P,Q$ 的类型相反，且 $(P,Q)$ 处于\Term{transversal-position}{横截位置}（参见 4.2.3）。
(iv) $P,Q$ 的类型相反，且 $\operatorname{rad}^{u}(P)\cap Q=e$。
(v) $\operatorname{rad}^{u}(P)\cap Q=\operatorname{rad}^{u}(Q)\cap P=e$。
(vi) \Term{canonical}{典范}\Term{morphism}{态射} $\operatorname{rad}^{u}(P)\times_SQ\to G$ 为\Term{open-immersion}{开浸入}。
(vi$'$) \Term{canonical}{典范}\Term{morphism}{态射} $\operatorname{rad}^{u}(P)\to G/Q$ 为\Term{open-immersion}{开浸入}。
(vii) 存在\Term{etale}{平展}\Term{topology}{拓扑}的\Term{covering-family}{覆盖族} $\{S_i\to S\}$，且对每个 $i$，有 $G_{S_i}$ 的\Term{splitting}{分裂} $(T_i,M_i,R_i)$ 及\Term{subset}{子集} $R_i^{(1)}\subset R_i$，使 $P_{S_i}$（或 $Q_{S_i}$）为 $G_{S_i}$ 中包含 $T_i$、由 $R_i^{(1)}$（或 $R_i^{(2)}=-R_i^{(1)}$）定义的 $(R)$ 型\Term{subgroup}{子群}。
\end{SGAresult}
\SourcePage{304}
\begin{proof}
先证第二部分。由 4.2.1(vii) 及\Term{splitting}{分裂}情形中 $s_G$ 的定义（第 XXII 讲，3.16.2(iv)），有 (iii)$\Leftrightarrow$(vii)；显然 (ii)$\Rightarrow$(i)；作\Term{base-change}{基变换} $G\to G/Q$ 得 (vi$'$)$\Rightarrow$(vi)。
现设 (vii) 成立，证明其余条件。因均关于\Term{etale}{平展}\Term{topology}{拓扑}\Term{local-scheme}{局部}，可假设 $G=(G,T,M,R)$ \Term{splitting}{分裂}，$P$ 由\Term{subset}{子集} $R'\subset R$ 定义，$Q$ 由 $-R'$ 定义。若 $L$ 为 $G$ 中包含 $T$、由 $R'\cap(-R')$ 定义的 $(R)$ 型\Term{subgroup}{子群}，则由第 XXII 讲，5.11.3，显然 $L$ 为 $P,Q$ 的共同\Term{levi-subgroup}{Levi 子群}。又 $P=L\cdot\operatorname{rad}^{u}(P)$、$Q=L\cdot\operatorname{rad}^{u}(Q)$，且由第 XXII 讲，5.6.7，$\operatorname{rad}^{u}(P)\cap Q=Q\cap\operatorname{rad}^{u}(P)=e$；故 $P\cap Q=L$，(ii)、(v) 已证。因 $P,Q$ 处于\Term{transversal-position}{横截位置}，\Term{canonical}{典范}\Term{morphism}{态射} $P\to G/Q$ 诱导\Term{open-immersion}{开浸入} $P/(P\cap Q)\to G/Q$（4.2.1）；而\Term{canonical}{典范}\Term{morphism}{态射} $\operatorname{rad}^{u}(P)\to P/(P\cap Q)=P/L$ 为\Term{isomorphism}{同构}，所以 (vi$'$) 已证。结合此前讨论，所有断言都是 (vii) 的推论。
现只须证明 (i)、(iv)、(v)、(vi) 中任意一项都蕴含 (vii)。既已证明 (ii)、(iii) \Term{equivalence}{等价}，只须在\Term{geometric-fibre}{几何纤维}上证明，可假设 $S$ 为\Term{algebraically-closed-field}{代数闭域}的谱。由 4.1.1，存在 $G$ 的\Term{maximal-torus}{极大环面} $T$ 包含在 $P\cap Q$ 中。令 $R$（或 $R_1,R_2$）为 $G$（或 $P,Q$）关于 $T$ 的\Term{root-set}{根集}。
令 $R_1^a$ 为 $R_1$ 的非对称部分（即 $R_1^a=\{\alpha\in R_1\mid-\alpha\notin R_1\}$），同样引入 $R_2^a$。须证 $R_1=-R_2$。
由 (i)，$R_1\cap R_2$ 对称。设 $\alpha\in R_1$；若 $\alpha\notin R_2$，则 $\alpha\in-R_2$；若 $\alpha\in R_2$，则 $\alpha\in R_1\cap R_2=-(R_1\cap R_2)\subset-R_2$。故 $R_1\subset-R_2$，由对称性 $R_1=-R_2$。
由 (iv)，$\operatorname{Card}(R_1)=\operatorname{Card}(R_2)$ 且 $R_1^a\cap R_2=\varnothing$；第二个条件\Term{equivalence}{等价}于 $R_2\subset-R_1$，再由第一个条件得 $R_2=-R_1$。
由 (v)，$R_1^a\cap R_2=R_2^a\cap R_1=\varnothing$，故 $R_2\subset-R_1$、$R_1\subset-R_2$，仍得 $R_2=-R_1$。
由 (vi)，$\operatorname{Lie}(\operatorname{rad}^{u}(P))\oplus\operatorname{Lie}(Q)=\operatorname{Lie}(G)$，故 $R$ 为 $R_2$ 与 $R_1^a$ 的\Term{disjoint-union}{不交并}，从而 $R_2=-R_1$。
第二部分证毕。再证第一部分。由 (vii)$\Rightarrow$(ii)，已证所求群 $P'$ 关于\Term{etale}{平展}\Term{topology}{拓扑}\Term{local-scheme}{局部}存在；仅须证明唯一性，这同样可关于\Term{etale}{平展}\Term{topology}{拓扑}\Term{local-scheme}{局部}进行。故可假设 $G$ 关于 $L$ 的\Term{maximal-torus}{极大环面} $T$ \Term{splitting}{分裂}，$P$（或 $P'$）由\Term{root-system}{根系} $R$ 的\Term{subset}{子集} $R_1$（或 $R'_1$）定义。
按假设 $R_1\cap R'_1$ 对称；如上论证，得 $R'_1=-R_1$。故 $P'$ 由 $P,L$ 决定，证明完成。
\end{proof}
\begin{SGAresult}{定义}{4.3.3}
满足 4.3.2 的\Term{equivalence}{等价}条件 (i) 至 (vii) 的 $G$ 的两个\Term{parabolic-subgroup}{抛物子群}称为\Term{opposite-parabolic}{相反的}。若 $P$ 为 $G$ 的\Term{parabolic-subgroup}{抛物子群}，$L$ 为 $P$ 的\Term{levi-subgroup}{Levi 子群}（或 $T$ 为 $P$ 的\Term{maximal-torus}{极大环面}），
\SourcePage{305}
则 $G$ 中使 $P\cap Q=L$ 的唯一\Term{parabolic-subgroup}{抛物子群} $Q$（或使 $P\cap Q$ 为 $P$ 中包含 $T$ 的唯一\Term{levi-subgroup}{Levi 子群}的 $Q$，参见 1.6；也即使 $P\cap Q$ 包含 $T$ 且 $P,Q$ \Term{opposite-parabolic}{相反}的 $Q$），称为相对于 $L$（或 $T$）与 $P$ \Term{opposite-parabolic}{相反}的\Term{parabolic-subgroup}{抛物子群}。
\end{SGAresult}
由 4.3.2(iii)，立即可从 4.2.4、4.2.5 得到平行的结果；举例如下。
\begin{SGAresult}{推论}{4.3.4}
设 $S$ 为\Term{scheme}{概形}，$G$ 为\Term{reductive-group}{约化} $S$-群，$P,Q$ 为 $G$ 的两个\Term{parabolic-subgroup}{抛物子群}。
(i) $P,Q$ \Term{opposite-parabolic}{相反}，当且仅当对每点 $s\in S$，$P_{\bar s},Q_{\bar s}$ \Term{opposite-parabolic}{相反}。若 $S'\to S$ \Term{surjectivity}{满射}且 $P_{S'},Q_{S'}$ \Term{opposite-parabolic}{相反}，则 $P,Q$ \Term{opposite-parabolic}{相反}。
(ii) \Term{functor}{函子} $\operatorname{Opp}(G)$（其中对 $S'\to S$，$\operatorname{Opp}(G)(S')$ 是 $G_{S'}$ 中\Term{opposite-parabolic}{相反}\Term{parabolic-subgroup}{抛物子群}对的\Term{set}{集合}）可由 $\operatorname{Par}(G)^2$ 的开子\Term{scheme}{概形}表示。\Term{functor}{函子} $\operatorname{Opp}(/P)$（其中 $\operatorname{Opp}(/P)(S')$ 是 $G_{S'}$ 中与 $P_{S'}$ \Term{opposite-parabolic}{相反}的\Term{parabolic-subgroup}{抛物子群}集）可由 $\operatorname{Par}_{s(t(P))}(G)$ 的一个在 $S$ 上\Term{universally-schematically-dense}{普遍概形稠密}的开子\Term{scheme}{概形}表示。\footnote{编者注：将“相对\Term{dense}{稠密}”改为“在 $S$ 上\Term{universally-schematically-dense}{普遍概形稠密}”，参见 EGA IV$_3$，定义 11.10.8。}
(iii) 设 $P,Q$ \Term{opposite-parabolic}{相反}；设 $P',Q'$ 为 $G$ 的两个\Term{parabolic-subgroup}{抛物子群}，$P'$ 与 $P$ 类型相同。则 $P',Q'$ \Term{opposite-parabolic}{相反}，当且仅当 $(P',Q')$ 关于\Term{etale}{平展}\Term{topology}{拓扑}\Term{local-scheme}{局部}与 $(P,Q)$ \Term{conjugacy}{共轭}。（注意：第 5 节将说明可将\Term{etale}{平展}\Term{topology}{拓扑}换成 Zariski \Term{topology}{拓扑}。）
\end{SGAresult}
\begin{SGAresult}{推论}{4.3.5}
设 $S$ 为\Term{scheme}{概形}，$G$ 为\Term{reductive-group}{约化} $S$-群，$P$ 为 $G$ 的\Term{parabolic-subgroup}{抛物子群}。
(i) \Term{set}{集合}上由 $Q\mapsto P\cap Q$ 定义的\Term{morphism}{态射} $\operatorname{Opp}(/P)\to\operatorname{Lev}(P)$（参见 1.9）为\Term{isomorphism}{同构}；$\operatorname{Opp}(/P)$ 是以 $\operatorname{rad}^{u}(P)$ 为结构群的\Term{torsor}{主齐性丛}（$\operatorname{rad}^{u}(P)$ 以\Term{inner-automorphism}{内自同构}\Term{action}{作用}）。若 $S$ \Term{affine}{仿射}，则存在 $G$ 的\Term{parabolic-subgroup}{抛物子群}与 $P$ \Term{opposite-parabolic}{相反}。
(ii) 设 $S$ 半\Term{local-scheme}{局部}，$\{s_i\}$ 为其\Term{closed-point}{闭点}集；对每个 $i$，设 $Q_i$ 为 $G_{s_i}$ 中与 $P_{s_i}$ \Term{opposite-parabolic}{相反}的\Term{parabolic-subgroup}{抛物子群}。则存在 $G$ 中与 $P$ \Term{opposite-parabolic}{相反}的\Term{parabolic-subgroup}{抛物子群} $Q$，使对每个 $i$ 有 $Q_{s_i}=Q_i$。
(iii) \Term{set}{集合}上由 $(P,Q)\mapsto(P,P\cap Q)$ 定义的\Term{morphism}{态射} $\operatorname{Opp}(G)\to\mathrm{PL}$（参见 3.15）为\Term{isomorphism}{同构}。
\end{SGAresult}
以上均由定理的第一部分以及 1.9、2.3、2.8 推出。
\textbf{注记 4.3.6。} 设 $P,Q$ 为 $G$ 的两个\Term{opposite-parabolic}{相反}\Term{parabolic-subgroup}{抛物子群}，$P\cdot Q$ 为 4.2.6 引入的 $G$ 的开子\Term{scheme}{概形}，即 $P\times_SQ$ 的层像。
\SourcePage{306}
“乘积\Term{morphism}{态射}”$G\times_SG\to G$ \Term{induction-isomorphism}{诱导同构}
\[\operatorname{rad}^{u}(P)\times_SQ\xrightarrow{\sim}P\cdot Q\xleftarrow{\sim}P\times_S\operatorname{rad}^{u}(Q).\]
这由 4.3.2（或 4.2.5(ii)）及 $P\cap Q$ 是 $P,Q$ 的共同\Term{levi-subgroup}{Levi 子群}立即推出，后者给出 $P=\operatorname{rad}^{u}(P)\cdot(P\cap Q)$、$Q=\operatorname{rad}^{u}(Q)\cdot(P\cap Q)$。
同样有\Term{commutative-diagram}{交换图}
\[\begin{tikzcd}
\operatorname{rad}^{u}(P)\arrow[r,hook]\arrow[d,"\sim"']&G/Q\arrow[d,"\sim"]\\
\operatorname{Opp}(/P)\arrow[r,hook]&\operatorname{Par}_{t(Q)}(G),
\end{tikzcd}\]
其中竖\Term{arrow}{箭头}由 $g\mapsto\operatorname{int}(g)Q$ 诱导。
\textbf{4.4. \Term{osculating-position}{密切位置}}
\begin{SGAresult}{命题}{4.4.1}
设 $S$ 为\Term{scheme}{概形}，$G$ 为\Term{reductive-group}{约化} $S$-群，$P,Q$ 为 $G$ 的两个\Term{parabolic-subgroup}{抛物子群}。下列条件\Term{equivalence}{等价}：
(i) $P\cap Q$ 是 $G$ 的\Term{parabolic-subgroup}{抛物子群}。
(ii) $P\cap Q$ 关于\Term{etale}{平展}\Term{topology}{拓扑}\Term{local-scheme}{局部}包含 $G$ 的\Term{borel-subgroup}{Borel 子群}。
(iii) $P\cap Q$ 关于\Term{etale}{平展}\Term{topology}{拓扑}\Term{local-scheme}{局部}包含 $G$ 的\Term{maximal-torus}{极大环面}，且对每个 $S'\to S$ 及 $G_{S'}$ 中包含于 $P_{S'},Q_{S'}$ 的每个\Term{maximal-torus}{极大环面} $T$，$P_{S'}$ 关于 $T$ 的\Term{opposite-parabolic}{相反子群}相对于 $Q_{S'}$ 处于\Term{transversal-position}{横截位置}。
(iv) 存在\Term{etale}{平展}\Term{topology}{拓扑}的\Term{covering-family}{覆盖族} $\{S_i\to S\}$，且对每个 $i$，有 $G_{S_i}$ 的\Term{maximal-torus}{极大环面} $T_i$ 包含于 $P_{S_i},Q_{S_i}$，并使 $P_{S_i}$ 关于 $T_i$ 的\Term{opposite-parabolic}{相反子群}相对于 $Q_{S_i}$ 处于\Term{transversal-position}{横截位置}。
(v) 存在\Term{etale}{平展}\Term{topology}{拓扑}的\Term{covering-family}{覆盖族} $\{S_i\to S\}$，且对每个 $i$，有 $G_{S_i}$ 的\Term{splitting}{分裂} $(T_i,M_i,R_i)$，使 $P_{S_i}$（或 $Q_{S_i}$）为 $G_{S_i}$ 中包含 $T_i$、由\Term{root-set}{根集} $R_i^{(1)}$（或 $R_i^{(2)}$）定义的 $(R)$ 型\Term{subgroup}{子群}，且 $R_i^{(1)}\cap R_i^{(2)}$ 包含 $R$ 的一个\Term{positive-root-system}{正根系}。
此外，若这些条件成立，则 $t(P\cap Q)=t(P)\cap t(Q)$（记号见 3.2）。
\end{SGAresult}
\begin{proof}
(v)$\Rightarrow$(ii)、(iii)$\Rightarrow$(iv) 显然。由 1.18，(ii)$\Rightarrow$(i)。有 (iv)$\Rightarrow$(v)：可假设 $G$ \Term{splitting}{分裂}，$P$（或 $Q$）由\Term{root-set}{根集} $R_1$（或 $R_2$）定义；$P$ 的\Term{opposite-parabolic}{相反子群}由 $-R_1$ 定义，故归结于引理 4.2.2。同样通过\Term{splitting}{分裂}证明 (i)$\Rightarrow$(iii)。最后，定理的末项可关于\Term{etale}{平展}\Term{topology}{拓扑}\Term{local-scheme}{局部}证明；可假设 $P\cap Q$ 包含 $G$ 的\Term{borel-subgroup}{Borel 子群} $B$，于是归结于 3.7。
\end{proof}
\begin{SGAresult}{定义}{4.4.2}
满足 4.4.1 的条件 (i) 至 (v) 的 $G$ 的两个\Term{parabolic-subgroup}{抛物子群}称为处于\Term{osculating-position}{密切位置（position osculatrice）}。
\end{SGAresult}
\SourcePage{307}
\begin{SGAresult}{推论}{4.4.3}
设 $P,Q$ 是处于\Term{osculating-position}{密切位置}的两个\Term{parabolic-subgroup}{抛物子群}，$P',Q'$ 为 $G$ 的两个\Term{parabolic-subgroup}{抛物子群}，分别与 $P,Q$ 类型相同。则 $P',Q'$ 处于\Term{osculating-position}{密切位置}，当且仅当 $(P',Q')$ 关于\Term{etale}{平展}\Term{topology}{拓扑}\Term{local-scheme}{局部}与 $(P,Q)$ \Term{conjugacy}{共轭}。
\end{SGAresult}
\begin{proof}
只须证明：若 $P,P'$ 均相对于 $Q$ 处于\Term{osculating-position}{密切位置}，则它们关于 $(\text{ét})$ \Term{local-scheme}{局部}由 $Q$ 的\Term{section}{截面}\Term{conjugacy}{共轭}。$P\cap Q,P'\cap Q$ 是包含于 $Q$ 的同类型的两个\Term{parabolic-subgroup}{抛物子群}，由下列引理 (ii)，它们关于 $(\text{ét})$ \Term{local-scheme}{局部}由 $Q$ 的\Term{section}{截面}\Term{conjugacy}{共轭}。故可假设 $P\cap Q=P'\cap Q$；此时由同一引理 (i)，$P=P'$。
\end{proof}
\begin{SGAresult}{引理}{4.4.4}
设 $P,P',Q$ 是\Term{reductive-group}{约化} $S$-群 $G$ 的三个\Term{parabolic-subgroup}{抛物子群}。
(i) $P=P'$，当且仅当 $P,P'$ 处于\Term{osculating-position}{密切位置}且类型相同。
(ii) 若 $P\subset Q$、$P'\subset Q$，且 $g\in G(S)$ 使 $\operatorname{int}(g)P$ 与 $P'$ 处于\Term{osculating-position}{密切位置}，则 $g\in Q(S)$。
\end{SGAresult}
\begin{proof}
(i) 由 4.4.1 的末项直接推出。证明 (ii)：$Q,\operatorname{int}(g)Q$ 均包含 $P'\cap\operatorname{int}(g)P$，故处于\Term{osculating-position}{密切位置}；由 (i) 它们相等，故 $g\in\operatorname{Norm}_G(Q)(S)=Q(S)$。
\end{proof}
注意，定理的 (iii)、(iv) 还立即给出：
\begin{SGAresult}{推论}{4.4.5}
设 $P,P',Q$ 是\Term{reductive-group}{约化} $S$-群 $G$ 的三个\Term{parabolic-subgroup}{抛物子群}，包含 $G$ 的同一\Term{maximal-torus}{极大环面} $T$。设 $P,P'$ 关于 $T$ \Term{opposite-parabolic}{相反}。则 $Q$ 相对于 $P$ 处于\Term{osculating-position}{密切位置}，当且仅当它相对于 $P'$ 处于\Term{transversal-position}{横截位置}。在这些条件下，$P\cap Q$ 也相对于 $P'$ 处于\Term{transversal-position}{横截位置}。
\end{SGAresult}
\begin{SGAresult}{推论}{4.4.6}
设 $P,Q$ 为 $G$ 的两个\Term{parabolic-subgroup}{抛物子群}，包含同一\Term{maximal-torus}{极大环面} $T$。则 $P,Q$ 处于\Term{transversal-position}{横截位置}，当且仅当存在 $G$ 的两个关于 $T$ \Term{opposite-parabolic}{相反}的\Term{parabolic-subgroup}{抛物子群} $P'\subset P$、$Q'\subset Q$；甚至可选取 $t(P')=t(P)\cap s(t(Q))$。\footnote{编者注：回顾，对合 $s$ 已在 4.3.1 定义。}
\end{SGAresult}
\begin{proof}
充分性显然（4.2.1(i)、4.3.2(iii)）。证明必要性。令 $P^-$（或 $Q^-$）为 $P$（或 $Q$）关于 $T$ 的\Term{opposite-parabolic}{相反子群}。由 4.4.5，$P^-\cap Q$ 相对于 $P,Q^-$ 处于\Term{transversal-position}{横截位置}，再用 4.4.5，它也相对于 $P\cap Q^-$ 处于\Term{transversal-position}{横截位置}；此外
\[\begin{aligned}
t(P^-\cap Q)&=t(P^-)\cap t(Q)=s(t(P))\cap s(t(Q^-))\\
&=s(t(P)\cap t(Q^-))=s(t(P\cap Q^-)).
\end{aligned}\]
故 $P^-\cap Q=P'$ 与 $P\cap Q^-=Q'$ \Term{opposite-parabolic}{相反}（4.3.2(iii)）；又 $P'\cap Q'\supset T$，故确实关于 $T$ \Term{opposite-parabolic}{相反}。
\end{proof}
\SourcePage{308}
\textbf{4.5. \Term{standard-position}{标准位置}}
本小节简要说明以上某些结果如何推广。
\begin{SGAresult}{命题}{4.5.1}
若 $P_1,P_2$ 是\Term{reductive-group}{约化} $S$-群 $G$ 的两个\Term{parabolic-subgroup}{抛物子群}，则下列条件\Term{equivalence}{等价}：\footnote{编者注：增添了这些条件\Term{equivalence}{等价}的证明（以及原著 6.17 隐含使用的条件 (v)）；相应地，将原 4.5.1 改为命题 4.5.1 及定义 4.5.1.1。}
(i) $P_1\cap P_2$ \Term{smooth}{光滑}。
(ii) $P_1\cap P_2$ 是 $G$ 的 $(R)$ 型（或 $(\mathrm{RC})$ 型）\Term{subgroup}{子群}。
(iii) $P_1\cap P_2$ 关于 $(\mathrm{fpqc})$ \Term{topology}{拓扑}\Term{local-scheme}{局部}包含 $G$ 的\Term{maximal-torus}{极大环面}。
(iv) $P_1\cap P_2$ 关于 Zariski \Term{topology}{拓扑}\Term{local-scheme}{局部}包含 $G$ 的\Term{maximal-torus}{极大环面}。
当 $S$ 半\Term{local-scheme}{局部}时，这些条件还\Term{equivalence}{等价}于：
(v) $P_1\cap P_2$ 包含 $G$ 的\Term{maximal-torus}{极大环面}。
\end{SGAresult}
\begin{proof}
显然 (iv)$\Rightarrow$(iii)（且当 $S$ 半\Term{local-scheme}{局部}时 (v)$\Rightarrow$(iv)），由第 XXII 讲，定义 5.2.1，(ii)$\Rightarrow$(i)。将证明 (i)$\Rightarrow$(ii)$\Rightarrow$(iii)，再证明 (iii) 蕴含 (i)、(iv)（且当 $S$ 半\Term{local-scheme}{局部}时也蕴含 (v)）。置 $K=P_1\cap P_2$。由 4.1.1 及 [BT65]，4.5，每个\Term{geometric-fibre}{几何纤维} $K_{\bar s}$ \Term{connected}{连通}，且包含 $G_{\bar s}$ 的\Term{maximal-torus}{极大环面} $T_{\bar s}$；$K_{\bar s}$ 关于 $T_{\bar s}$ 的\Term{root-set}{根集}是 $G_{\bar s}$ 关于 $T_{\bar s}$ 的\Term{root-set}{根集}的\Term{closed-subset}{闭子集}。因此，若 $K$ 在 $S$ 上\Term{smooth}{光滑}，它就是 $(\mathrm{RC})$ 型（从而也是 $(R)$ 型）\Term{subgroup}{子群}，参见第 XXII 讲，5.2.1、5.11.1。故 (i)$\Leftrightarrow$(ii)。
若 $K$ 为 $(R)$ 型\Term{subgroup}{子群}，由第 XXII 讲，2.2 及第 XIX 讲，6.1，它关于\Term{etale}{平展}\Term{topology}{拓扑}\Term{local-scheme}{局部}包含 $G$ 的\Term{maximal-torus}{极大环面}。故 (ii)$\Rightarrow$(iii)。
设 (iii) 成立，证明 $K$ \Term{smooth}{光滑}。由 $(\mathrm{fpqc})$ \Term{descent}{下降}，可假设 $G=(G,T,M,R)$ \Term{splitting}{分裂}，其中 $T$ 为包含于 $P\cap Q$ 的\Term{maximal-torus}{极大环面}，且有 $R$ 的两个\Term{closed-subset}{闭子集} $R_1,R_2$ 使 $P=H_{R_1}$、$Q=H_{R_2}$。因 $K$ \Term{fiber}{纤维}\Term{connected}{连通}，第 XXII 讲，5.4.5 给出 $K=H_{R_1\cap R_2}$，故它在 $S$ 上\Term{smooth}{光滑}且为 $(\mathrm{RC})$ 型。
令 $R_1^s$ 为 $R_1$ 的对称部分，$R_1^a=R_1-R_1^s$；则 $\operatorname{rad}^{u}(P_1)=H_{R_1^a}$。如 [BT65]，4.4 的证明中所指出，$R'=(R_1^s\cap R_2)\cup R_1^a$ 是 $R$ 的\Term{closed-subset}{闭子集}，且 $R'\cup(-R')=R$；故由第 XXI 讲，3.3.6，$R'$ 包含 $R$ 的\Term{positive-root-system}{正根系}。此外，$R'$ 的对称部分是 $R_1^s\cap R_2^s$，包含于 $R_1\cap R_2$。
由以上可知，$P'=K\cdot\operatorname{rad}^{u}(P)$ 是 $G$ 的\Term{parabolic-subgroup}{抛物子群}，且 $K$ 是 $P'$ 的 $(\mathrm{RC})$ 型\Term{subgroup}{子群}，满足 $\operatorname{rad}^{u}(K)=K\cap\operatorname{rad}^{u}(P')$。由 2.11，$K$ 有\Term{levi-subgroup}{Levi 子群} $L$。再由第 XIV 讲，3.20、3.21，$K$ 满足 (iv)，当 $S$ 半\Term{local-scheme}{局部}时还满足 (v)。4.5.1 证毕。
\end{proof}
\SourcePage{309}
\begin{SGAresult}{定义}{4.5.1.1}
\footnote{编者注：参见脚注 (22)。关于不处于\Term{standard-position}{标准位置}的\Term{parabolic-subgroup}{抛物子群} $P,Q$ 的例子，参见下文 7.11。}若以上条件成立，则称 $P,Q$ 相互处于\Term{standard-position}{标准位置}。例如，若 $P,Q$ 处于\Term{transversal-position}{横截位置}、或\Term{osculating-position}{密切位置}，或\Term{base}{基底}为域的谱，便是如此。此概念在\Term{base-extension}{基扩张}下\Term{stable}{稳定}，且关于 $(\mathrm{fpqc})$ \Term{topology}{拓扑}\Term{local-scheme}{局部}。
\end{SGAresult}
\textbf{4.5.2.} 设 $(P,Q),(P',Q')$ 为 $G$ 的两对处于\Term{standard-position}{标准位置}的\Term{parabolic-subgroup}{抛物子群}，令 $H$ 为 $G$ 的下列\Term{subfunctor}{子函子}：$H(S')$ 是满足 $\operatorname{int}(g)P=P'$ 且 $\operatorname{int}(g)Q=Q'$ 的 $g\in G(S')$ 所成的\Term{set}{集合}。它是 $G$ 的闭子\Term{scheme}{概形}，在 $S$ 上\Term{smooth}{光滑}，且\Term{form}{形式}上是以 $P\cap Q$ 为结构群的\Term{torsor}{主齐性丛}。由此得到下列条件\Term{equivalence}{等价}：
(i) $(P,Q),(P',Q')$ 关于 $(\mathrm{fpqc})$ \Term{topology}{拓扑}\Term{local-scheme}{局部}\Term{conjugacy}{共轭}；
(ii) $(P,Q),(P',Q')$ 关于\Term{etale}{平展}\Term{topology}{拓扑}\Term{local-scheme}{局部}\Term{conjugacy}{共轭}；
(iii) $(P,Q),(P',Q')$ 在每条\Term{geometric-fibre}{几何纤维}上\Term{conjugacy}{共轭}。
此时称 $(P,Q),(P',Q')$ 有相同的相互位置类型。这一概念在\Term{base-change}{基变换}下\Term{stable}{稳定}，且关于 $(\mathrm{fpqc})$ \Term{topology}{拓扑}\Term{local-scheme}{局部}。
\textbf{4.5.3.} 令 $\operatorname{Stand}(G)$ 为 $\operatorname{Par}(G)\times_S\operatorname{Par}(G)$ 中“由相互处于\Term{standard-position}{标准位置}的对构成”的\Term{subfunctor}{子函子}。则 $\operatorname{Stand}(G)$ \Term{representable-functor}{可表示}，存在\Term{finite}{有限}\Term{etale}{平展} $S$-\Term{scheme}{概形} $\operatorname{TypeStand}$（“相互\Term{standard-position}{标准位置}类型\Term{scheme}{概形}”），以及\Term{smooth}{光滑}、\Term{finite-presentation}{有限表现}、\Term{geometric-fibre}{几何纤维}\Term{irreducible}{不可约}（特别地\Term{faithfully-flat}{忠实平坦}）的\Term{morphism}{态射}
\[t_2:\operatorname{Stand}(G)\longrightarrow\operatorname{TypeStand},\]
它是 $\operatorname{Stand}(G)$ 对 $G$ \Term{action}{作用}的商：$\operatorname{Stand}(G)$ 的两个\Term{section}{截面}（在任意 $S'\to S$ 上）具有相同相互位置类型，当且仅当它们在 $t_2$ 下的像相同。有\Term{commutative-diagram}{交换图}
\[\begin{tikzcd}[column sep=large]
\operatorname{Stand}(G)\arrow[r,"t_2"]\arrow[d,hook]&\operatorname{TypeStand}\arrow[d,"q"]\\
\operatorname{Par}(G)\times_S\operatorname{Par}(G)\arrow[r,"t\times t"]&\operatorname{Of}(\operatorname{Dyn}(G))\times_S\operatorname{Of}(\operatorname{Dyn}(G)),
\end{tikzcd}\]
其中\Term{morphism}{态射} $q$ 可经\Term{descent}{下降}如下描述：若 $(P,Q)$ 是 $G$ 的一对相对处于\Term{standard-position}{标准位置}的\Term{parabolic-subgroup}{抛物子群}，且 $T$ 为 $P\cap Q$ 的\Term{maximal-torus}{极大环面}，则\Term{set}{集合}上由 $n\mapsto(P,\operatorname{int}(n)Q)$ 定义的\Term{morphism}{态射} $\operatorname{Norm}_G(T)\to\operatorname{Stand}(G)$ \Term{induction-isomorphism}{诱导同构}
\[\mathrm W_P(T)\backslash\mathrm W_G(T)/\mathrm W_Q(T)\simeq q^{-1}(t(P),t(Q)).\]
（左端表示\Term{double-coset}{双陪集}\Term{sheaf}{层}……。）这些断言不难证明（特别注意 $t_2^{-1}(t_2(P,Q))\simeq G/(P\cap Q)$）。
\SourcePage{310}
\textbf{4.5.4.} 现固定 $G$ 的一个\Term{parabolic-subgroup}{抛物子群} $P$，以 $\operatorname{Par}(G;P)$ 表示 $G$ 中相对于 $P$ 处于\Term{standard-position}{标准位置}的\Term{parabolic-subgroup}{抛物子群}\Term{functor}{函子}。对每个 $t\in\operatorname{Of}(\operatorname{Dyn}(G))(S)$，同样置 $\operatorname{Par}_t(G;P)=\operatorname{Par}(G;P)\cap\operatorname{Par}_t(G)$。立即可见，这两个\Term{functor}{函子}由 $\operatorname{Stand}(G)$ 作\Term{fiber-product}{纤维积}得到，故可由在 $S$ 上\Term{smooth}{光滑}、\Term{finite-presentation}{有限表现}、\Term{fiber}{纤维}\Term{nonempty}{非空}的 $S$-\Term{scheme}{概形}表示。$t_2$ 诱导\Term{canonical}{典范}\Term{morphism}{态射} $t_P$（即 $t_P(Q)=t_2(P,Q)$）
\[t_P:\operatorname{Par}_t(G;P)\longrightarrow q^{-1}(t(P),t),\]
它\Term{smooth}{光滑}、\Term{finite-presentation}{有限表现}、\Term{geometric-fibre}{几何纤维}\Term{irreducible}{不可约}。\Term{canonical}{典范}\Term{morphism}{态射} $\operatorname{Par}_t(G;P)\to\operatorname{Par}_t(G)$ 是\Term{surjectivity}{满射}\Term{monomorphism}{单态射}，故可视为 $\operatorname{Par}_t(G)$ 的一个\Term{cellular-decomposition}{胞腔分解}（由 $q^{-1}(t(P),t)$ 的\Term{connected-component}{连通分支}集编号）。
\textbf{4.5.5.} 现设类型 $t$ 形如 $t(Q)$，其中 $Q$ 为 $G$ 中相对于 $P$ 处于\Term{standard-position}{标准位置}的\Term{parabolic-subgroup}{抛物子群}，且 $P\cap Q$ 包含\Term{maximal-torus}{极大环面} $T$。
此时 $\operatorname{Par}_t(G)\simeq G/Q$，$q^{-1}(t(P),t)\simeq\mathrm W_P(T)\backslash\mathrm W_G(T)/\mathrm W_Q(T)$，得到图
\[\begin{tikzcd}[column sep=large]
\operatorname{Par}_{t(Q)}(G;P)\arrow[r,"f"]\arrow[d,hook,"i"']&\mathrm W_P(T)\backslash\mathrm W_G(T)/\mathrm W_Q(T)\\
\operatorname{Par}_{t(Q)}(G)\arrow[r,"\sim"]&G/Q,
\end{tikzcd}\]
其中 $i$ 为\Term{surjectivity}{满射}\Term{monomorphism}{单态射}，$f$ \Term{smooth}{光滑}、\Term{finite-presentation}{有限表现}、\Term{geometric-fibre}{几何纤维}\Term{irreducible}{不可约}。此外，若 $Q_1,Q_2$ 是 $\operatorname{Par}_{t(Q)}(G;P)$ 的两个\Term{section}{截面}（在 $S'\to S$ 上），即为 $G_{S'}$ 中与 $Q$ 关于 $(\mathrm{fpqc})$ \Term{local-scheme}{局部}\Term{conjugacy}{共轭}、且相对于 $P_{S'}$ 处于\Term{standard-position}{标准位置}的两个\Term{parabolic-subgroup}{抛物子群}，则 $Q_1,Q_2$ 关于 $(\mathrm{fpqc})$ \Term{local-scheme}{局部}由 $P$ 的\Term{section}{截面}\Term{conjugacy}{共轭}，当且仅当 $f(Q_1)=f(Q_2)$。当 $S$ 为\Term{algebraically-closed-field}{代数闭域} $k$ 的谱时，得到关系
\[P(k)\backslash G(k)/Q(k)\simeq\mathrm W_P(T)(k)\backslash\mathrm W_G(T)(k)/\mathrm W_Q(T)(k).\]
一般地，若假设\Term{scheme}{概形} $\mathrm W_P(T)\backslash\mathrm W_G(T)/\mathrm W_Q(T)$ 为常值且形如 $E_S$（例如 $G$ 关于 $T$ \Term{splitting}{分裂}且 $S$ \Term{connected}{连通}时即如此），则 $f^{-1}(e)$（$e\in E$）构成 $\operatorname{Par}_{t(Q)}(G;P)$ 到开闭子\Term{scheme}{概形}的\Term{resolution}{分解}；这些子\Term{scheme}{概形}为 $P$ 的\Term{homogeneous-space}{齐性空间}，在 $S$ 上\Term{smooth}{光滑}、\Term{finite-presentation}{有限表现}、\Term{geometric-fibre}{几何纤维}\Term{irreducible}{不可约}。
\textbf{4.5.6.} 回到 4.5.4 的一般情形。\Term{scheme}{概形} $q^{-1}(t(P),t)$\footnote{编者注：将 $t^{-1}(t(P),t)$ 纠正为 $q^{-1}(t(P),t)$，并在下文两处将 $\operatorname{Par}_t(G)$ 纠正为 $\operatorname{Par}_t(G;P)$。}总有两个特殊\Term{section}{截面}，分别对应“\Term{transversal-position}{横截位置}”及“\Term{osculating-position}{密切位置}”两种类型。第一个\Term{section}{截面}的\Term{inverse-image}{逆像}
\SourcePage{311}
如上所述，是 $\operatorname{Par}_t(G;P)$ 的相对\Term{dense-open-set}{稠密开集}，即此\Term{resolution}{分解}中\Term{relative-dimension}{相对维数}最大的胞腔。
第二个\Term{section}{截面}的\Term{inverse-image}{逆像}大概是 $\operatorname{Par}_t(G;P)$ 的闭子\Term{scheme}{概形}；它是此\Term{resolution}{分解}中\Term{relative-dimension}{相对维数}最小的胞腔。
